\chapter{Currents, positivity, and analytic cycles}
\index{current!CLD}\index{positivity!current}\index{analytic cycle}
\label{chap:cld-currents}

Currents are continuous linear functionals on compactly supported smooth
forms.  This apparently formal definition has three geometric consequences:
singular curvature becomes meaningful, analytic cycles acquire intrinsic
integration currents, and weak convergence becomes available.  The price is
that products of arbitrary currents do not exist.  We shall keep every
defined product visible in the notation.

Let $X$ be a topologically separated, pure $n$-dimensional analytic space.
Unless a boundary term is explicitly retained, $X$ is assumed boundaryless.

\section{Orientation: what a current is for}
\label{sec:ch14-orientation}

The three preceding chapters built a sheaf of smooth forms
$A_X^{p,q}$ on a Berkovich space, an integral
$\int_X:A_{X,c}^{n,n}\to\mathbf R$ normalized by lattice measures and
analytic multiplicities, and a Stokes formula
(\Cref{cor:analytic-stokes-green}).  Every object so far has been
\emph{smooth}: a form is, by construction, the pullback of a superform with
$C^\infty$ coefficients along a tropical chart.  This is already enough to
formulate curvature for a smooth metric, but it is not enough for anything
singular.  Three demands force the enlargement.

\begin{enumerate}
\item \emph{Singular curvature.}  The function $u=\max\{0,t\}$ on the
  tropical line is not smooth, yet it is the local model for the potential
  of a model metric on a curve.  Its second derivative exists only as a
  measure, and one wants to write $\dprimeddprime u=\delta_0$.
\item \emph{Cycles.}  A closed analytic subspace $Y\subset X$ of dimension
  $r<n$ carries no smooth $(n-r,n-r)$-form on $X$ representing it, but
  it does define a linear functional
  $\alpha\mapsto\int_Yi^*\alpha$ on $A_{X,c}^{r,r}$.  Poincar\'e--Lelong
  (\Cref{chap:cld-poincare-lelong}) is an identity between such
  functionals.
\item \emph{Limits.}  Envelopes, monotone limits of psh functions, and
  Bedford--Taylor products are all constructed by passing to a limit which is
  not uniform in any $C^k$ norm.  The only topology in which such limits
  converge is the weak one.
\end{enumerate}

Each demand is met by the same device: replace the object by the functional
it induces on test forms.  The cost is that a functional cannot in general be
multiplied by another functional, and this chapter is written so that every
product it does define is visible in the notation.  Where a product is only
defined under hypotheses, those hypotheses are displayed in the statement
rather than deferred to a remark.

\subsection{What is proved here, and what is imported}
\label{sec:ch14-prerequisites}

The reader needs the superform algebra of \Cref{chap:superforms}---in
particular the three positivity cones of \Cref{def:ch09-cones}, their matrix
description \Cref{prop:ch09-positivity-dictionary}, the duality theorem
\Cref{thm:ch09-positivity-duality}, and the product rules
\Cref{prop:ch09-positivity-products}---together with the calibrated
integration theory of \Cref{chap:tropical-integration} and the analytic
sheaves of \Cref{chap:cld-forms}.  From \Cref{chap:cld-forms} we use above
all the description of the test-form topology,
\Cref{prop:test-form-continuity-criterion}, and the existence of smooth
partitions of unity, \Cref{thm:cld-smooth-partition}.  From general topology
we use Tychonoff's theorem and the Riesz representation theorem; the latter
is recalled in \Cref{chap:topological-tools}.

Everything in
\Crefrange{sec:ch14-definition}{sec:ch14-dictionary} is proved except for the
following imported inputs, each of which is used with an exact citation at the
point of use:
\begin{enumerate}
\item the continuity and piecewise smoothness of fibre integrals in the
  general, non-skeletal case, \cite[Theorem~16.3.3]{CLD25}, which is what
  licenses the use of $u_\omega$ as the coefficient of a test form in the
  limit argument of \Cref{thm:cld-fibre-projection};
\item Lagerberg's Euclidean vanishing theorem for positive closed
  supercurrents, \cite[Corollary~4.1]{Lag12Super}, used through
  \cite[Theorem~16.8.5]{CLD25};
\item Berndtsson's construction of convex Bedford--Taylor products,
  \cite[Theorem~4.2]{Ber26Convex}, and his minimality criterion,
  \cite[Corollary~5.2]{Ber19Super}, both quoted verbatim in
  \Cref{sec:ch14-berndtsson,sec:ch14-convex-bt}.
\end{enumerate}
No other statement of this chapter is used without proof.

\subsection{The plan of the chapter}
\label{sec:ch14-plan}

By the end of the chapter the reader should be able to
\begin{enumerate}
\item write down the continuity condition defining a current, and produce a
  current whose order is not bounded over chart data;
\item pass between the three positivity cones for forms and the three
  positivity cones for currents, and know which of the two orderings is
  reversed;
\item compute $\dprime T$, $\ddprime T$, and $\dprimeddprime T$ with the
  correct signs, including the boundary term;
\item push a current forward along a proper morphism, and say precisely why
  pullback is unavailable;
\item recognize the integration current of an analytic cycle and of a
  skeleton, and compute both on the Berkovich line;
\item carry out a fibre integration to a number, including the case where a
  degree appears;
\item state exactly what fails, in the non-Archimedean setting, of the
  complex compactness and regularization package, and exhibit the example
  that fails;
\item distinguish Berndtsson's metric supercurrent of a real submanifold from
  the calibrated tropical current of the same support, and compute the
  discrepancy.
\end{enumerate}
The final section collects a dictionary, a table of competing
normalizations, and the unit tests.

\section{Definition, bidegree, and weak topology}
\label{sec:ch14-definition}

\subsection{The topology on test forms, restated}
\label{sec:ch14-test-topology}

Because every later statement about weak convergence refers to it, we restate
the topology of \Cref{chap:cld-forms} in the exact form in which it will be
used.  Fix $p,q\ge0$.

\begin{definition}[Chart datum]
\label{def:ch14-chart-datum}
\index{chart datum}
A \emph{chart datum} on $X$ is a pair $(\mathscr D,K)$ consisting of a finite
family
\[
 \mathscr D=\bigl(f_i:U_i\to\mathbf G_m^{r_i}\bigr)_{i=1}^N
\]
of moment maps on open subsets of $X$ and a compact set
$K\subset\bigcup_iU_i$.  The letter $K$ is used for the compact set
throughout this chapter whenever a chart datum is in play; the ground field,
elsewhere denoted $K$, will be written explicitly when it is needed.  We write $A_{\mathscr D,K}^{p,q}(X)$ for the space of
smooth $(p,q)$-forms $\omega$ with $\Supp\omega\subset K$ that admit, on each
$U_i$, a representation $\omega|_{U_i}=f_{i,\trop}^*\alpha_i$.
\end{definition}

For a compact subset $W$ of some $U_i$ and an integer $r\ge0$ one has the
seminorm $\|\cdot\|_{W,r}$ of \Cref{chap:cld-forms}: choose a finite
polyhedral decomposition of $f_{i,\trop}(W)$ and affine coordinates on each
cell, and take the maximum of the absolute values of all partial derivatives
of order $\le r$ of all coefficients of $\alpha_i$ on the relevant cells.  Two
such choices give equivalent seminorms.  The space
$A_{\mathscr D,K}^{p,q}(X)$ carries the locally convex topology defined by all
these seminorms, and
\[
 A_{X,c}^{p,q}=\varinjlim_{(\mathscr D,K)}A_{\mathscr D,K}^{p,q}(X)
\]
carries the final locally convex topology.  \emph{Chart data are finite
families}: this is what makes \Cref{prop:test-form-continuity-criterion}
usable, and it is exactly the point at which the non-Archimedean theory
diverges from the classical one, since a finite family of charts on a compact
Berkovich space need not exhaust the local tropical directions available
there.  \Cref{ex:ch14-star-unbounded-order} below turns this remark into a
counterexample.

\begin{definition}
\label{def:ch14-current}
A \emph{current of bidimension $(r,s)$} is a continuous linear functional
\[
                     T:A_{X,c}^{r,s}\longrightarrow\mathbf R.
\]
The space of such currents is denoted $D_{r,s}(X)$.  We put
\[
                    D^{p,q}(X)=D_{n-p,n-q}(X)
\]
and call its elements currents of bidegree $(p,q)$.  The \emph{weak
topology} is pointwise convergence on test forms:
\[
 T_i\longrightarrow T
 \quad\Longleftrightarrow\quad
 \langle T_i,\alpha\rangle\longrightarrow\langle T,\alpha\rangle
 \quad\text{for every }\alpha\in A_{X,c}^{n-p,n-q}.
\]
\end{definition}

Continuity refers exactly to the final $C^\infty$ topology constructed in
\Cref{chap:cld-forms}.  A current has order at most $r$ on fixed
compact chart data if it satisfies the corresponding $C^r$ estimate.  There
need not be a single order valid for every chart datum.

\begin{definition}[Relative order]
\label{def:ch14-relative-order}
\index{current!order of}
Let $T\in D_{r,s}(X)$ and let $(\mathscr D,K)$ be a chart datum.  The
\emph{order of $T$ relative to $(\mathscr D,K)$} is the least integer
$m\ge0$, if one exists, for which there are a constant $C>0$ and finitely many
compact chart domains $W_1,\ldots,W_k$ with
\begin{equation}
\label{eq:ch14-order-bound}
 |\langle T,\omega\rangle|\le C\max_j\|\omega\|_{W_j,m}
 \qquad\text{for all }\omega\in A_{\mathscr D,K}^{r,s}(X).
\end{equation}
We say that $T$ has \emph{global order $\le m$} if the same $m$ works for
every chart datum, with a constant $C$ allowed to depend on the compact set
but not on the chart datum; equivalently, if $T$ is continuous for the $C^m$
topology on $A_{X,c}^{r,s}$.
\end{definition}

By \Cref{prop:test-form-continuity-criterion} a linear functional is a current
if and only if it has a finite relative order for each chart datum.  The
distinction between relative and global order is not pedantry: it is the
hinge on which \Cref{sec:ch14-compactness} turns.

\begin{example}[Forms with locally integrable coefficients]
\label{ex:ch14-l1-current}
Let $\omega$ be a $G$-form of bidegree $(p,q)$ on $X$ whose coefficients are
tropically Borel and locally bounded, or more generally locally integrable
for the measures attached to the charts.  Then
\[
 \langle[\omega],\alpha\rangle=\int_X\omega\wedge\alpha,
 \qquad \alpha\in A_{X,c}^{n-p,n-q},
\]
defines a current of bidegree $(p,q)$ of global order $0$: on a chart datum
$(\mathscr D,K)$ the integral is bounded by
$\|\omega\|_{L^1(K)}\|\alpha\|_{K,0}$ up to the calibration weights, which are
bounded on $K$.  The assignment $\omega\mapsto[\omega]$ is linear and weakly
continuous, and it is compatible with the notation
$[\mu]$ for a measure $\mu$ used below.  Its injectivity is a separate matter
and is discussed in \Cref{rem:ch14-form-injectivity}.
\end{example}

\begin{example}[Forms and measures]
\label{ex:ch14-forms-and-measures}
If $\omega$ is a locally integrable $(p,q)$-form, then
\[
               [\omega](\alpha)=\int_X\omega\wedge\alpha
\]
defines a current of bidegree $(p,q)$.  A Radon measure $\mu$ defines a
current of bidimension $(0,0)$ by
$\langle[\mu],f\rangle=\int f\,\rd \mu$.  Conversely, positive currents of
bidimension $(0,0)$ are precisely positive Radon measures.
\end{example}

\begin{proof}[Proof of the last assertion]
On a compact set choose a nonnegative smooth cutoff $\chi$ equal to one.
For a smooth test function $f$ supported there,
$-\|f\|_\infty\chi\leq f\leq\|f\|_\infty\chi$.  Positivity gives
\[
                  |T(f)|\leq T(\chi)\|f\|_\infty.
\]
Thus $T$ extends continuously from smooth to continuous compactly supported
functions, using smooth Stone--Weierstrass density, and the Riesz
representation theorem supplies the Radon measure.  Uniqueness follows from
the same density.
\end{proof}

The existence of the cutoff $\chi$ is \Cref{thm:cld-smooth-partition}, and the
density of smooth compactly supported functions in $C_c(X)$ is the
Stone--Weierstrass step inside its proof; both need $X$ locally
holomorphically separated and paracompact, which we assume from here on
whenever a partition of unity is used.  Observe that the argument used only
the \emph{positivity} of the linear functional, not its continuity.  That this
is a general phenomenon is the content of
\Cref{prop:ch14-automatic-continuity} below.

\subsection{Currents form a fine sheaf}
\label{sec:ch14-sheaf}

For an open $U\subset X$ let $j:U\hookrightarrow X$ be the inclusion.
Extension by zero of compactly supported test forms is a continuous linear map
$j_!:A_{U,c}^{r,s}\to A_{X,c}^{r,s}$, since a chart datum on $U$ is one on
$X$.  Its transpose is the \emph{restriction} $T\mapsto T|_U$,
\[
 \langle T|_U,\alpha\rangle=\langle T,j_!\alpha\rangle,
 \qquad \alpha\in A_{U,c}^{r,s}.
\]

\begin{proposition}
\label{prop:ch14-currents-sheaf}
Let $X$ be locally holomorphically separated and paracompact.  Then
$U\mapsto D_{r,s}(U)$ is a sheaf of $A_X^{0,0}$-modules, and for each integer
$m$ the currents of global order $\le m$ form a subsheaf.  All these sheaves
are fine.  In particular the order of a current is a local notion.
\end{proposition}

\begin{proof}
\emph{Presheaf.}  Restriction is transitive because extension by zero is, and
it is linear; so $U\mapsto D_{r,s}(U)$ is a presheaf.  The $A^{0,0}$-module
structure is $\langle uT,\alpha\rangle=\langle T,u\alpha\rangle$ for a smooth
function $u$, which is compatible with restriction.

\emph{Separatedness.}  Let $(U_i)$ cover $U$ and let $T\in D_{r,s}(U)$ satisfy
$T|_{U_i}=0$ for all $i$.  Let $\alpha\in A_{U,c}^{r,s}$.  By
\Cref{thm:cld-smooth-partition} there are finitely many smooth compactly
supported functions $\varphi_i$, with $\Supp\varphi_i\subset U_i$, such that
$\sum_i\varphi_i=1$ on a neighbourhood of $\Supp\alpha$; only finitely many
are needed because $\Supp\alpha$ is compact.  Then
$\alpha=\sum_i\varphi_i\alpha$ and each $\varphi_i\alpha$ is compactly
supported in $U_i$, so $\langle T,\alpha\rangle=\sum_i\langle
T|_{U_i},\varphi_i\alpha\rangle=0$.  Hence $T=0$.

\emph{Gluing.}  Let $T_i\in D_{r,s}(U_i)$ agree on overlaps.  For
$\alpha\in A_{U,c}^{r,s}$ choose $(\varphi_i)$ as above and set
\[
 \langle T,\alpha\rangle=\sum_i\langle T_i,\varphi_i\alpha\rangle .
\]
\emph{Independence of $(\varphi_i)$.}  If $(\varphi'_j)$ is a second such
family, then $\sum_{i,j}\varphi'_j\varphi_i=1$ near $\Supp\alpha$, and
$\varphi'_j\varphi_i\alpha$ is compactly supported in $U_i\cap U_j$, so
\[
 \sum_j\langle T_j,\varphi'_j\alpha\rangle
 =\sum_{i,j}\langle T_j,\varphi'_j\varphi_i\alpha\rangle
 =\sum_{i,j}\langle T_i,\varphi'_j\varphi_i\alpha\rangle
 =\sum_i\langle T_i,\varphi_i\alpha\rangle,
\]
the middle equality being the agreement of $T_i$ and $T_j$ on
$U_i\cap U_j$.  Thus $T$ is well defined, and it is visibly linear.

\emph{Continuity.}  Fix a chart datum $(\mathscr D,K)$ on $U$.  Choose
$(\varphi_i)$ equal to $1$ near $K$, with $\Supp\varphi_i\subset U_i$; this is
a finite family.  For $\omega\in A^{r,s}_{\mathscr D,K}(U)$ we have
$\langle T,\omega\rangle=\sum_i\langle T_i,\varphi_i\omega\rangle$, and
$\omega\mapsto\varphi_i\omega$ is continuous from $A^{r,s}_{\mathscr D,K}(U)$
into $A_{U_i,c}^{r,s}$ because multiplication by a fixed smooth function is
continuous for each seminorm by the Leibniz rule.  A finite sum of continuous
functionals is continuous, so $T$ is a current.

\emph{Restriction.}  If $\Supp\alpha$ is compact in $U_j$, then
$\Supp(\varphi_i\alpha)\subset U_i\cap U_j$, so
$\langle T_i,\varphi_i\alpha\rangle=\langle T_j,\varphi_i\alpha\rangle$ and
$\langle T,\alpha\rangle=\langle T_j,\alpha\rangle$.  Thus $T|_{U_j}=T_j$.

\emph{Order.}  In the continuity estimate just written, if each $T_i$ has
order $\le m$ with constant $C_i$ on the relevant compacta, then $T$ has order
$\le m$ with constant $\sum_iC_i\max_i\|\varphi_i\|_{C^m}$.  Hence
currents of order $\le m$ form a subsheaf.  Conversely a current of order
$\le m$ restricts to one of order $\le m$.

\emph{Fineness.}  Multiplication by a locally finite smooth partition of unity
subordinate to a cover gives sheaf endomorphisms with locally finite supports
summing to the identity, which is the definition of a fine sheaf.  The
partition exists by \Cref{thm:cld-smooth-partition}.
\end{proof}

\begin{definition}[Support]
\label{def:ch14-support}
\index{current!support of}
The support $\Supp(T)$ of a current is the complement of the largest open set
on which $T$ vanishes.  It is well defined by
\Cref{prop:ch14-currents-sheaf}, and $\langle T,\alpha\rangle=0$ whenever
$\Supp\alpha$ is compact and disjoint from $\Supp T$.
\end{definition}

\begin{lemma}
\label{lem:ch14-phi-support}
Let $T\in D_{r,s}(X)$ and let $\alpha\in A_X^{r,s}$ be a smooth form, not
necessarily compactly supported, such that $\Supp\alpha\cap\Supp T$ is
compact.  Then the number
\[
 \langle T,\alpha\rangle:=\langle T,\lambda\alpha\rangle
\]
is independent of the choice of a smooth compactly supported $\lambda$ equal
to $1$ near $\Supp\alpha\cap\Supp T$.  The resulting extension of $T$ is
linear and continuous for the topology of uniform convergence of all
derivatives on the family of closed sets meeting $\Supp T$ compactly.  In
particular, if $\Supp T$ is compact, then $\langle T,\alpha\rangle$ is defined
for every smooth $\alpha$.
\end{lemma}

\begin{proof}
Such a $\lambda$ exists by \Cref{thm:cld-smooth-partition}.  If $\lambda'$ is a
second choice, then $\lambda-\lambda'$ vanishes near
$\Supp\alpha\cap\Supp T$, so
$\Supp((\lambda-\lambda')\alpha)\cap\Supp T=\emptyset$ and
$\langle T,(\lambda-\lambda')\alpha\rangle=0$ by \Cref{def:ch14-support}.
Homogeneity is clear.  For additivity, given $\alpha,\beta$ choose one
$\lambda$ equal to $1$ near
$(\Supp\alpha\cup\Supp\beta)\cap\Supp T$ and use linearity of $T$.  For
continuity, fix a closed $F$ with $F\cap\Supp T$ compact and choose $\lambda$
equal to $1$ near $F\cap\Supp T$; then
$\langle T,\alpha\rangle=\langle T,\lambda\alpha\rangle$ for every $\alpha$
supported in $F$, and $\alpha\mapsto\lambda\alpha$ is continuous into
$A_{X,c}^{r,s}$.
\end{proof}

\subsection{Extension by zero and removable sets}
\label{sec:ch14-extension}

The following lemma is the first place at which the non-Archimedean theory
departs sharply from the complex one.  Its hypothesis is a statement about
\emph{tropical} dimension, and it will be seen in
\Cref{sec:ch14-compactness} to make many sets removable that are not removable
over $\mathbf C$.

\begin{lemma}
\label{lem:ch14-extension-by-zero}
Let $U\subset X$ be open, $j:U\hookrightarrow X$ the inclusion, and
$T\in D_{r,s}(U)$.  Suppose that the local tropical dimension satisfies
\[
 d_{\trop}(x)<\max\{r,s\}\qquad\text{for every }x\in X\setminus U .
\]
Then every smooth compactly supported $(r,s)$-form on $X$ has support
contained in $U$, the map $\alpha\mapsto\alpha|_U$ is an isomorphism of
topological vector spaces $A_{X,c}^{r,s}\to A_{U,c}^{r,s}$, and consequently
restriction
\[
 D_{r,s}(X)\longrightarrow D_{r,s}(U)
\]
is an isomorphism, with inverse $T\mapsto j_!T$ defined by
$\langle j_!T,\alpha\rangle=\langle T,\alpha|_U\rangle$.
\end{lemma}

\begin{proof}
Let $\alpha\in A_{X,c}^{r,s}$ and let $x\in X\setminus U$.  By
\Cref{prop:analytic-tropical-dimension} the local tropical dimension at $x$
is computed on a tropical chart, and by the dimension vanishing recalled in
\Cref{chap:cld-forms} (see the remark following
\Cref{prop:cld-form-pullback}) a superform of bidegree $(r,s)$ restricted to a
polyhedron of dimension $<\max\{r,s\}$ vanishes, because one of its two
alternating factors has degree exceeding the dimension.  Hence $\alpha$
vanishes identically on a neighbourhood of $x$.  As $x$ was arbitrary,
$\Supp\alpha\subset U$.

Therefore restriction $A_{X,c}^{r,s}\to A_{U,c}^{r,s}$ is a bijection with
inverse $j_!$; both directions are continuous because a chart datum on one
side is a chart datum on the other with the same seminorms.  Dualizing gives
the assertion about currents; the formula for $j_!T$ is the transpose of the
inverse map.
\end{proof}

\begin{corollary}
\label{cor:ch14-rigid-invisible}
Let $X$ be a good $K$-analytic space of pure dimension $n$, and let
$Z\subset X$ be a closed weakly analytic subset with
$\dim Z<\max\{r,s\}$.  Then restriction
$D_{r,s}(X)\to D_{r,s}(X\setminus Z)$ is an isomorphism; equivalently, no
nonzero current of bidimension $(r,s)$ is supported in $Z$.  In particular,
on the analytification of a smooth projective curve, no nonzero current of
bidimension $(1,1)$ is supported on the set of rigid points.
\end{corollary}

\begin{proof}
A weakly analytic subset of dimension $d$ has tropical dimension $\le d$ at
each of its points by \Cref{prop:analytic-tropical-dimension}, so
\Cref{lem:ch14-extension-by-zero} applies with $U=X\setminus Z$.  For the last
sentence one must not treat the full rigid locus as a closed weakly analytic
subset; it need not be one.  Instead, if a current $T$ of
bidimension $(1,1)$ were supported on the rigid points, take
$U=X\setminus\Supp T$.  The complement is closed, every one of its points is
rigid, and hence it has local tropical dimension $0$ pointwise.  Thus
\Cref{lem:ch14-extension-by-zero} makes restriction to $U$ injective; since
$T|_U=0$, one has $T=0$.
\end{proof}

\begin{warning}{Removability is not a positivity theorem}
\Cref{lem:ch14-extension-by-zero} is a statement about \emph{all} currents of
a given bidimension, and its hypothesis is a dimension inequality that fails
as soon as $\max\{r,s\}=0$.  It says nothing about bidimension $(0,0)$
currents, that is, about measures: a Dirac mass at a rigid point is a
perfectly good positive current of bidimension $(0,0)$.  Nor does it say that
a current on $X\setminus Z$ of bidimension $(0,0)$ extends.  Extension
theorems in the borderline case require growth or positivity hypotheses; see
\Cref{chap:cld-poincare-lelong,chap:cld-bedford-taylor}.
\end{warning}

\section{Differentials and their signs}
\label{sec:ch14-differentials}

Let $T\in D^{p,q}(X)$.  The current differentials are defined by
\[
\begin{aligned}
 \langle\dprime T,\alpha\rangle
   &=(-1)^{p+q+1}\langle T,\dprime\alpha\rangle,\\
 \langle\ddprime T,\alpha\rangle
   &=(-1)^{p+q+1}\langle T,\ddprime\alpha\rangle.
\end{aligned}
\]
The exponent uses the \emph{bidegree} of $T$ before applying the operator.

\subsection{Where the sign comes from}
\label{sec:ch14-sign-derivation}

The sign is not a convention that could have been chosen otherwise; it is
forced by the requirement that the inclusion of smooth forms into currents be
a morphism of complexes.  We derive it once, in the form that will be reused
in \Cref{chap:cld-poincare-lelong,chap:delta-forms}.

Write $b'$ for the naked transpose of $\dprime$, that is
$\langle b'T,\alpha\rangle=\langle T,\dprime\alpha\rangle$, and similarly
$b''$.  These are the maps that exist for purely formal reasons:
$\dprime:A_{X,c}^{r,s}\to A_{X,c}^{r+1,s}$ is continuous by
\Cref{prop:cld-intrinsic-calculus}, so its transpose maps
$D_{r+1,s}(X)\to D_{r,s}(X)$, and it is weakly continuous because the weak
topology is the topology of pointwise convergence.  From $\ddprime=J\dprime
J^{-1}$ on forms one gets $b''=Jb'J^{-1}$ on currents, and from
$(\dprime)^2=(\ddprime)^2=0$, $\dprime\ddprime=-\ddprime\dprime$ one gets
$(b')^2=(b'')^2=0$ and $b'b''=-b''b'$.

\begin{lemma}
\label{lem:ch14-sign-forced}
Let $\omega$ be a smooth $(p,q)$-form on a boundaryless pure
$n$-dimensional $X$, and let $\alpha\in A_{X,c}^{n-p-1,n-q}$.  Then
\[
 \langle b'[\omega],\alpha\rangle
 =(-1)^{p+q+1}\langle[\dprime\omega],\alpha\rangle .
\]
Consequently the unique constant $\varepsilon(p,q)\in\{\pm1\}$ for which
$T\mapsto\varepsilon(p,q)b'T$ satisfies $\dprime[\omega]=[\dprime\omega]$ for
all smooth $\omega$ of bidegree $(p,q)$ is $\varepsilon(p,q)=(-1)^{p+q+1}$.
\end{lemma}

\begin{proof}
By the graded Leibniz rule of \Cref{prop:cld-intrinsic-calculus},
\[
 \dprime(\omega\wedge\alpha)
 =\dprime\omega\wedge\alpha+(-1)^{p+q}\omega\wedge\dprime\alpha .
\]
The form $\omega\wedge\alpha$ has bidegree $(n-1,n)$ and compact support, so
$\int_X\dprime(\omega\wedge\alpha)=0$ by
\Cref{cor:analytic-stokes-green}.  Integrating the displayed identity
therefore gives
\[
 0=\int_X\dprime\omega\wedge\alpha
   +(-1)^{p+q}\int_X\omega\wedge\dprime\alpha
 =\langle[\dprime\omega],\alpha\rangle
   +(-1)^{p+q}\langle b'[\omega],\alpha\rangle,
\]
whence $\langle b'[\omega],\alpha\rangle=(-1)^{p+q+1}
\langle[\dprime\omega],\alpha\rangle$, because $(-1)^{p+q}$ is its own
inverse.  Uniqueness of $\varepsilon(p,q)$ is clear since
$\omega\mapsto[\omega]$ is not identically zero in each bidegree.
\end{proof}

Accordingly we set, once and for all,
\begin{equation}
\label{eq:ch14-dprime-sign}
 \dprime=(-1)^{p+q+1}b'
 \quad\text{and}\quad
 \ddprime=(-1)^{p+q+1}b''
 \qquad\text{on }D^{p,q}(X),
\end{equation}
which is the definition displayed at the head of the section.  This is the
convention recorded in \Cref{sec:appJ-currents}; the reader who prefers to
remember a test rather than a sign should remember
\Cref{lem:ch14-sign-forced} and the boundary formula of
\Cref{prop:ch14-boundary-current} below.

\begin{proposition}
\label{prop:current-differential-calculus}
The operators on currents are continuous for the weak topology and satisfy
\[
 (\dprime)^2=(\ddprime)^2=0,
 \qquad \dprime\ddprime+\ddprime\dprime=0,
 \qquad \ddprime=J\dprime J^{-1}.
\]
Moreover,
\[
             \langle\dprimeddprime T,\alpha\rangle
              =\langle T,\dprimeddprime\alpha\rangle.
\]
If $X$ is boundaryless and $\omega$ is smooth, then
$\dprime[\omega]=[\dprime\omega]$ and
$\ddprime[\omega]=[\ddprime\omega]$.
\end{proposition}
\begin{proof}
Transpose the identities on test forms, keeping the displayed signs.  The
two signs in the mixed second derivative multiply to $-1$, which is canceled
by $\ddprime\dprime=-\dprime\ddprime$ on test forms; this gives the last
pairing formula.  For a smooth form, integration by parts and
\Cref{cor:analytic-stokes-green} identify the distributional and
smooth differentials.  With boundary, the omitted Stokes term is generally
nonzero.

We give the computations in full.  \emph{Continuity} holds because the weak
topology is that of pointwise convergence and $\dprime$, $\ddprime$ are
defined by composing with a fixed continuous operator on test forms.

\emph{Squares.}  On $D^{p,q}$ we have $\dprime=(-1)^{p+q+1}b'$ and, on
$D^{p+1,q}$, $\dprime=(-1)^{p+q}b'$; hence on $D^{p,q}$,
\[
 \dprime\dprime=(-1)^{p+q}(-1)^{p+q+1}(b')^2=-(b')^2=0 .
\]
The same computation gives $(\ddprime)^2=0$.

\emph{Anticommutation.}  Starting from $T\in D^{p,q}$, the composite
$\dprime\ddprime$ picks up $(-1)^{p+q+1}$ from $\ddprime$ and
$(-1)^{p+q+2}$ from the subsequent $\dprime$ acting on $D^{p,q+1}$, a total
sign of $-1$; therefore $\dprime\ddprime=-b'b''$.  Symmetrically
$\ddprime\dprime=-b''b'$.  Since $b'b''=-b''b'$, we get
$\dprime\ddprime=-\ddprime\dprime$ as claimed.

\emph{Conjugation by $J$.}  If $T\in D^{p,q}$, then $JT\in D^{q,p}$ has the
same total degree, so the two operators carry the same sign and
$\ddprime=J\dprime J^{-1}$ follows from $b''=Jb'J^{-1}$.

\emph{The $\dprimeddprime$ pairing.}  By the preceding paragraph
$\dprimeddprime T=-b'b''T$, and $-b'b''$ is the transpose of
$-\ddprime\dprime=\dprimeddprime$ on test forms.  Hence
$\langle\dprimeddprime T,\alpha\rangle=\langle T,\dprimeddprime\alpha\rangle$
with no sign at all.  This is the identity that makes
$\dprimeddprime$ symmetric and is the reason it, rather than $\dprime$ or
$\ddprime$ separately, is the operator of potential theory.

\emph{Smooth forms.}  This is \Cref{lem:ch14-sign-forced} for $\dprime$; for
$\ddprime$ apply $J$ and use $J[\omega]=[J\omega]$, which holds because
$\int_XJ\eta=\int_X\eta$ for a top-degree form, the involution acting
trivially on $\tau_n$ by \eqref{eq:appJ-lagerberg-orientation}.
\end{proof}

\begin{proposition}
\label{prop:ch14-boundary-current}
Let $X$ be pure $n$-dimensional, possibly with analytic boundary, and let
$\omega$ be a $G$-form of bidegree $(p,q)$ with tropically $C^1$
coefficients.  Denote by $[\omega]^\partial$ the current
$\alpha\mapsto\int_X^\partial\omega\wedge\alpha$ of bidimension
$(n-p-1,n-q)$ built from the boundary integral of
\Cref{chap:tropical-integration}.  Then
\[
 \dprime[\omega]=[\dprime\omega]-[\omega]^\partial,
 \qquad
 \ddprime[\omega]=[\ddprime\omega]-J\bigl[J^{-1}\omega\bigr]^\partial .
\]
In particular, if $X$ is boundaryless and $\omega$ is a form (rather than a
$G$-form), both boundary terms vanish and
$\dprime[\omega]=[\dprime\omega]$, $\ddprime[\omega]=[\ddprime\omega]$.
\end{proposition}

\begin{proof}
Repeat the computation of \Cref{lem:ch14-sign-forced}, but use the Stokes
formula with boundary term,
$\int_X\dprime\eta=\int_X^\partial\eta$, of
\Cref{thm:cld-global-stokes} instead of
\Cref{cor:analytic-stokes-green}.  With
$\eta=\omega\wedge\alpha$ this gives
\[
 \int_X^\partial\omega\wedge\alpha
 =\int_X\dprime\omega\wedge\alpha
   +(-1)^{p+q}\int_X\omega\wedge\dprime\alpha,
\]
that is $[\omega]^\partial=[\dprime\omega]-\dprime[\omega]$ after multiplying
by $(-1)^{p+q}$ and using \eqref{eq:ch14-dprime-sign}.  For $\ddprime$ apply
$J$ to the identity already proved for $J^{-1}\omega$, using
$J\dprime J^{-1}=\ddprime$ and $J[\eta]=[J\eta]$.  If $X$ is boundaryless and
$\omega$ is an ordinary form, then the boundary measure attached to any test
form vanishes, so $[\omega]^\partial=0$.
\end{proof}

\begin{corollary}
\label{cor:ch14-delta-x-closed}
If $X$ is boundaryless, the integration current
$\delta_X=[1]$ of bidimension $(n,n)$ is $\dprime$- and $\ddprime$-closed.
\end{corollary}

\begin{proof}
Apply \Cref{prop:ch14-boundary-current} with $\omega=1$, a smooth
$(0,0)$-form: $\dprime1=0$ and the boundary term vanishes.
\end{proof}

\subsection{Multiplication by a smooth form}
\label{sec:ch14-wedge-smooth}

If $\alpha$ is smooth, set
\[
               \langle T\wedge\alpha,\beta\rangle
                =\langle T,\alpha\wedge\beta\rangle.
\]
This is always defined.  By contrast, $T\wedge S$ has no meaning for two
arbitrary currents.

\begin{lemma}
\label{lem:ch14-module-leibniz}
Let $T\in D_{r,s}(X)$ and let $\alpha\in A_X^{p,q}$, $\beta\in A_X^{p',q'}$.
\begin{enumerate}
\item $T\wedge\alpha$ is a current of bidimension $(r-p,s-q)$, of order
  $\le m$ if $T$ is; the map $T\mapsto T\wedge\alpha$ is weakly continuous and
  $(T\wedge\alpha)\wedge\beta=T\wedge(\alpha\wedge\beta)$.
\item $J(T\wedge\alpha)=JT\wedge J\alpha$, and setting
  $\alpha\wedge T=(-1)^{(p+q)(r+s)}T\wedge\alpha$ makes
  $D_{\bullet,\bullet}(X)$ a graded bimodule over
  $A_X^{\bullet,\bullet}$.
\item With $T$ of bidegree $(u,v)$ one has
  \[
   \dprime(T\wedge\alpha)=(\dprime T)\wedge\alpha
    +(-1)^{u+v}T\wedge\dprime\alpha .
  \]
\end{enumerate}
\end{lemma}

\begin{proof}
(1) and (2) are immediate from associativity and the multiplicativity of $J$
on forms, \Cref{prop:cld-intrinsic-calculus}; continuity of
$\beta\mapsto\alpha\wedge\beta$ on each $A^{r,s}_{\mathscr D,K}$ follows from
the Leibniz rule applied to the coefficients in a chart, which produces a
bound by the $C^m$ seminorms of $\beta$ times fixed $C^m$ seminorms of
$\alpha$.

(3) Let $T$ have bidegree $(u,v)$ and let $\beta$ be a test form of
complementary bidegree for the left-hand side.  Writing $w=u+v$ and
$w'=p+q$, and using \eqref{eq:ch14-dprime-sign} at the two relevant
bidegrees together with the Leibniz rule on forms,
\[
\begin{aligned}
 \langle\dprime(T\wedge\alpha),\beta\rangle
 &=(-1)^{w+w'+1}\langle T\wedge\alpha,\dprime\beta\rangle
  =(-1)^{w+w'+1}\langle T,\alpha\wedge\dprime\beta\rangle\\
 &=(-1)^{w+w'+1}(-1)^{w'}
   \bigl\langle T,\dprime(\alpha\wedge\beta)-\dprime\alpha\wedge\beta
   \bigr\rangle\\
 &=(-1)^{w+1}\langle T,\dprime(\alpha\wedge\beta)\rangle
   +(-1)^{w}\langle T,\dprime\alpha\wedge\beta\rangle\\
 &=\langle(\dprime T)\wedge\alpha,\beta\rangle
   +(-1)^{w}\langle T\wedge\dprime\alpha,\beta\rangle,
\end{aligned}
\]
which is the assertion.
\end{proof}

\section{Positivity cones for currents}
\label{sec:ch14-positivity}

The duality reverses the strength of positivity.  This is not a quirk of the
definition but a theorem about cones: by
\Cref{thm:ch09-positivity-duality} the dual of the strongly positive cone is
the weakly positive one and conversely, while the positive cone is
self-dual.  A current is a functional on forms, so testing it against the
\emph{largest} cone of forms is the \emph{strongest} requirement.  The
resulting hierarchy therefore reads, for currents,
\[
 \{\text{strongly positive}\}\subset\{\text{positive}\}
   \subset\{\text{weakly positive}\},
\]
exactly as for forms, but with the roles of the two extreme test cones
exchanged.  Getting this exchange wrong is the commonest error in the
subject, and \Cref{prop:ch14-form-current-cones} below is the statement that
makes it harmless: for currents defined by forms the two hierarchies agree
term by term.

\begin{definition}
A symmetric current $T$ of bidimension $(p,p)$ is
\begin{enumerate}
\item weakly positive if it is nonnegative on strongly positive test forms;
\item positive if it is nonnegative on positive test forms;
\item strongly positive if it is nonnegative on weakly positive test forms.
\end{enumerate}
Equivalently, a current of bidegree $(n-p,n-p)$ is tested in the same way.
\end{definition}

\begin{proposition}
\label{prop:current-positivity-calculus}
Positivity is local and weakly closed.  Proper pushforward preserves all
three positivity notions.  Wedge product with a smooth form has the following
properties:
\begin{enumerate}
\item if $\alpha$ is strongly positive, then $T\wedge\alpha$ has the same
      positivity type as $T$;
\item if $\alpha$ and $T$ are positive, then $T\wedge\alpha$ is positive;
\item if $\alpha$ is weakly positive and $T$ is strongly positive, then
      $T\wedge\alpha$ is weakly positive.
\end{enumerate}
Positive currents are of order zero on fixed compact chart data.
\end{proposition}
\begin{proof}
Locality follows by multiplying a positive test form by a nonnegative
partition of unity.  Weak closure follows by passing to the limit in every
nonnegative pairing.  Pullback preserves positivity of test forms, hence its
transpose, proper pushforward, preserves positivity of currents.  The wedge
statements follow by combining the corresponding wedge-product rules for
smooth forms with the reversed test cones in the definition of positivity for
currents.  In particular, the wedge of two positive smooth superforms is
positive.  Finally, on compact chart data a symmetric form is bounded above
and below, up to a fixed constant, by a fixed strongly positive form times its
$C^0$ norm.  Testing these two inequalities against $T$ gives an order-zero
bound.  The wedge-product assertions are
\cite[Lemma~5.2.9 and Proposition~3.2.4]{CLD25}.

We now supply the details, since each of the four assertions is used later.

\emph{Locality.}  Suppose $T|_{U_i}$ has the given positivity type for a cover
$(U_i)$ of $X$, and let $\alpha$ be a compactly supported test form in the
corresponding cone.  Choose smooth $\varphi_i\ge0$ with
$\Supp\varphi_i\subset U_i$ and $\sum_i\varphi_i=1$ near $\Supp\alpha$
(\Cref{thm:cld-smooth-partition}).  Each $\varphi_i\alpha$ lies in the same
cone, because all three cones are stable under multiplication by a
nonnegative function (\Cref{prop:ch09-positivity-inclusions}) and this is a
pointwise statement.  Hence
$\langle T,\alpha\rangle=\sum_i\langle T|_{U_i},\varphi_i\alpha\rangle\ge0$.
The converse, that restriction preserves positivity, follows because
extension by zero of a form in a cone stays in that cone.  Thus each
positivity type defines a subsheaf of $D_{p,p}$.

\emph{Weak closure.}  If $T_i\to T$ weakly and each $T_i$ is nonnegative on a
fixed cone of test forms, then so is $T$, the inequality
$\langle T_i,\alpha\rangle\ge0$ passing to the limit.  Symmetry passes to the
limit as well, since $J$ is weakly continuous.

\emph{Proper pushforward.}  If $g:Y\to X$ is topologically proper, then
$g^*$ preserves each cone of forms, by
\Cref{prop:cld-form-pullback} combined with
\Cref{prop:ch09-positivity-inclusions}, whose pullback statement is proved
there for affine maps and therefore holds chartwise.  Hence
$\langle g_*T,\alpha\rangle=\langle T,g^*\alpha\rangle\ge0$ whenever
$\alpha$ lies in the test cone attached to the positivity type of $T$.

\emph{Wedge products.}  Assertion (1): let $\alpha$ be strongly positive and
let $T$ be, say, weakly positive; for a strongly positive test form $\beta$
the product $\alpha\wedge\beta$ is strongly positive by
\Cref{prop:ch09-positivity-products}\,(1), so
$\langle T\wedge\alpha,\beta\rangle=\langle T,\alpha\wedge\beta\rangle\ge0$.
The same argument with \Cref{prop:ch09-positivity-products}\,(2) and (3)
gives the cases of $T$ positive and $T$ strongly positive.
Assertion (2): if $\alpha$ and $\beta$ are positive, so is
$\alpha\wedge\beta$ by \Cref{prop:ch09-positivity-products}\,(2), whence
$T\wedge\alpha$ is positive when $T$ is.  Assertion (3): if $\alpha$ is
weakly positive and $\beta$ is strongly positive, then $\alpha\wedge\beta$ is
weakly positive by \Cref{prop:ch09-positivity-products}\,(3); testing a
strongly positive $T$ against it gives $\ge0$, so $T\wedge\alpha$ is weakly
positive.  Notice that in each case the \emph{test} cone of the conclusion is
determined by which factor was allowed to be the weaker one; this is the
current-level shadow of the warning after
\Cref{prop:ch09-positivity-products}.

\emph{Order zero.}  Fix a chart datum $(\mathscr D,K)$ and a chart
$f:U\to\mathbf G_m^r$ of $\mathscr D$.  Give the target
$\mathbf R^r$ a Euclidean norm and put
$\eta=\dprimeddprime\|x\|^2$, a strongly positive $(1,1)$-superform with
constant coefficients.  Let $W\subset U$ be a compact chart domain
containing $\Supp\rho$ for a fixed cutoff $\rho$.  By
\Cref{lem:ch14-eta-dominates} below, $\eta^p$ is an interior point of the
positive cone $\mathrm{P}_p$; hence there is $c>0$ with
$\alpha\pm c\|\alpha\|_{W,0}\,\eta^{p}$ positive for every
symmetric $\alpha$ on $f_{\trop}(W)$.  Pulling back preserves positivity by
\Cref{prop:ch09-positivity-inclusions}, and multiplying by
$\rho\ge0$ preserves it as well, so testing against the positive current $T$
gives
\[
 |\langle T,\rho f_{\trop}^*\alpha\rangle|
 \le c\,\langle T,\rho f_{\trop}^*\eta^p\rangle\,\|\alpha\|_{W,0},
\]
an order-zero bound.  Since the constant $c\langle
T,\rho f_{\trop}^*\eta^p\rangle$ depends on the chart, this proves order zero
\emph{relative to each chart datum}, which is all that is claimed.
\end{proof}

\begin{lemma}
\label{lem:ch14-eta-dominates}
Let $V$ be a real vector space of dimension $r$ with a Euclidean norm, let
$\eta=\dprimeddprime\|x\|^2$, and let $0\le p\le r$.  Write $|\alpha|$ for
the operator norm of the normalized coefficient matrix $M(\alpha)$ of
\Cref{def:ch09-coefficient-matrix}.  Then
\[
 M(\eta^{p})=2^{p}p!\,\mathrm{Id},
\]
and for every symmetric $\alpha\in A^{p,p}(V)$ with constant coefficients and
every $c\ge(2^pp!)^{-1}$,
\[
 \alpha+c\,|\alpha|\,\eta^{p}\in\mathrm{P}_p(V)\subset\mathrm{WP}_p(V).
\]
Thus $\eta^p$ lies in the interior of the positive cone and in the interior
of the weakly positive cone.
\end{lemma}

\begin{proof}
In an orthonormal basis $x_1,\ldots,x_r$ the Hessian convention
\eqref{eq:appJ-real-hessian} gives
$\eta=2\sum_i\dprime x_i\wedge\ddprime x_i$.  Expanding the $p$-th power, the
terms with a repeated index vanish, and the $p!$ orderings of a fixed
$p$-element set contribute equally because each factor
$\dprime x_i\wedge\ddprime x_i$ has even total degree and therefore commutes
with everything, by \Cref{prop:ch09-wedge-signs}.  Hence
\[
 \eta^{p}=2^{p}p!\sum_{I\in S_p}
   \dprime x_{i_1}\wedge\ddprime x_{i_1}\wedge\cdots
   \wedge\dprime x_{i_p}\wedge\ddprime x_{i_p}
 =2^{p}p!\sum_{I\in S_p}c_p\,\dprime x_I\wedge\ddprime x_I,
\]
the sum being over increasing multi-indices $I=(i_1<\cdots<i_p)$ and the
second equality being \Cref{lem:ch09-generators}.  In the normalization of
\Cref{def:ch09-coefficient-matrix} this reads $M(\eta^p)=2^pp!\,\mathrm{Id}$.

By \Cref{prop:ch09-positivity-dictionary}\,(1) a symmetric $(p,p)$-form is
positive exactly when its matrix is positive semidefinite.  Since
$M(\alpha)\succeq-|\alpha|\,\mathrm{Id}$ by definition of the operator norm,
for $c\ge(2^pp!)^{-1}$ we get
\[
 M\bigl(\alpha+c|\alpha|\eta^p\bigr)
 =M(\alpha)+c\,2^pp!\,|\alpha|\,\mathrm{Id}
 \succeq M(\alpha)+|\alpha|\,\mathrm{Id}\succeq0 .
\]
The inclusion $\mathrm{P}_p\subset\mathrm{WP}_p$ is
\Cref{prop:ch09-positivity-inclusions}.  Finally $\mathrm{Id}$ is positive
definite, and both the condition ``$M\succeq0$'' of
\Cref{prop:ch09-positivity-dictionary}\,(1) and the condition
``$w^{\mathsf T}Mw\ge0$ for $w\in\Gamma_p$'' of
\Cref{prop:ch09-positivity-dictionary}\,(3) hold strictly at $\mathrm{Id}$ on
the compact set of unit vectors involved; so $\eta^p$ is an interior point of
both cones.
\end{proof}

\begin{warning}{$\eta^p$ is not interior to the strongly positive cone}
For $r\ge4$ and $2\le p\le r-2$ the cone $\mathrm{SP}_p(V)$ does not span the
space of symmetric $(p,p)$-forms, and therefore has empty interior in it.
Indeed the weakly null form $\gamma$ of \Cref{ex:ch14-weakly-null} satisfies
$\langle M(\gamma),vv^{\mathsf T}\rangle=v^{\mathsf T}M(\gamma)v=0$ for every
decomposable $v$, so $M(\gamma)\neq0$ is orthogonal to
$M(\mathrm{SP}_p)$ for the trace pairing.  What is true, and is
\cite[Proposition~1.2.13]{CLD25}, is that $\eta^p$ lies in the
\emph{relative} interior of $\mathrm{SP}_p(V)$, that is, in its interior
inside the linear span of the cone.  Consequently every argument in this book
that dominates an arbitrary symmetric form by a multiple of $\eta^p$
concludes positivity or weak positivity, never strong positivity.  This is
the precise reason why \Cref{prop:ch14-automatic-continuity} is stated for
functionals nonnegative on positive, or on weakly positive, forms.
\end{warning}

\begin{proposition}
\label{prop:ch14-form-current-cones}
Let $X$ be pure $n$-dimensional and boundaryless, and let $\omega$ be a
symmetric smooth $(p,p)$-form on $X$.  If $\omega$ is strongly positive
(respectively positive, weakly positive), then the current $[\omega]$ of
bidegree $(p,p)$, that is of bidimension $(n-p,n-p)$, is strongly positive
(respectively positive, weakly positive).
\end{proposition}

\begin{proof}
Suppose $\omega$ is strongly positive and let $\beta$ be a weakly positive
test form of bidegree $(n-p,n-p)$.  By
\Cref{thm:ch09-positivity-duality}, $\mathrm{WP}_{n-p}$ is the dual cone of
$\mathrm{SP}_p$, so $\omega\wedge\beta$ is a nonnegative multiple of
$\tau_n$ at every point of every chart; hence
$\int_X\omega\wedge\beta\ge0$, because the integral of a nonnegative
top-degree form against a positive calibration is nonnegative.  Thus
$[\omega]$ is strongly positive.  The positive case uses the self-duality
$\mathrm{P}_p=\mathrm{P}_{n-p}^\vee$, and the weakly positive case the
duality $\mathrm{WP}_p=\mathrm{SP}_{n-p}^\vee$.  Symmetry of $[\omega]$
follows from $J[\omega]=[J\omega]=[\omega]$.
\end{proof}

\begin{remark}[Injectivity and the converse]
\label{rem:ch14-form-injectivity}
The converse of \Cref{prop:ch14-form-current-cones}---that $[\omega]$
positive forces $\omega$ positive---requires knowing that a nonzero
continuous form is detected by a skeleton of full dimension, and this is a
genuine analytic input about the richness of the space rather than a formal
consequence of duality.  For a very rich tropical space, in particular for a
boundaryless Berkovich space of pure dimension $n$, it holds, and the map
$\omega\mapsto[\omega]$ is injective on forms with tropically continuous
coefficients; this is \cite[Corollaries~7.3.9 and~7.3.10]{CLD25}.  We use
only the direction proved above, except in
\Cref{sec:ch14-compactness}, where the injectivity is quoted with this
citation.
\end{remark}

\begin{proposition}
\label{prop:ch14-automatic-continuity}
Let $X$ be locally holomorphically separated and paracompact, let $p\ge0$,
and let $S:A_{X,c}^{p,p}\to\mathbf R$ be a linear map---\emph{not} assumed
continuous---such that $S(J\omega)=S(\omega)$ for all $\omega$.
\begin{enumerate}
\item If $S(\omega)\ge0$ for every positive $\omega$, then $S$ is a positive
  current, of order $0$ relative to every chart datum.
\item If $S(\omega)\ge0$ for every weakly positive $\omega$, then $S$ is a
  strongly positive current, of order $0$ relative to every chart datum.
\end{enumerate}
\end{proposition}

\begin{proof}
Only continuity has to be proved; symmetry is the hypothesis $S\circ J=S$,
and the positivity type of the conclusion is the hypothesis read through
\Cref{def:ch09-cones}.  We write the proof of (1); for (2) replace
$\mathrm{P}_p$ by $\mathrm{WP}_p$ throughout, which is legitimate because
$\eta^p$ is an interior point of both cones by
\Cref{lem:ch14-eta-dominates}.

By \Cref{prop:test-form-continuity-criterion} it suffices to bound $S$ on
each $A_{\mathscr D,K}^{p,p}(X)$, and by a partition of unity subordinate to
the finitely many chart domains of $\mathscr D$ it suffices to treat forms of
the shape $\rho\,f_{\trop}^*\alpha$, where $f:U\to\mathbf G_m^r$ is one
chart, $\rho\ge0$ is smooth with compact support in $U$, and $\alpha$ is a
symmetric superform on a compact polyhedral $W\supset f_{\trop}(\Supp\rho)$.

Let $\eta=\dprimeddprime\|x\|^2$ on $\mathbf R^r$ and let $c$ be the constant
of \Cref{lem:ch14-eta-dominates}.  Then
\[
 \alpha+c\|\alpha\|_{W,0}\,\eta^{p}
\]
is positive at every point of $W$, hence so is its pullback by
\Cref{prop:ch09-positivity-inclusions}, hence so is the product with
$\rho\ge0$, since multiplying by a nonnegative function preserves
$\mathrm{P}_p$ pointwise.  Applying $S$ and using its hypothesis,
\[
 S\bigl(\rho f_{\trop}^*\alpha\bigr)
  +c\|\alpha\|_{W,0}\,S\bigl(\rho f_{\trop}^*\eta^p\bigr)\ge0 .
\]
Applying the same inequality to $-\alpha$ and combining,
\[
 \bigl|S\bigl(\rho f_{\trop}^*\alpha\bigr)\bigr|
 \le c\,S\bigl(\rho f_{\trop}^*\eta^p\bigr)\,\|\alpha\|_{W,0},
\]
which is the required $C^0$ bound; the constant
$c\,S(\rho f_{\trop}^*\eta^p)$ is finite and nonnegative, being the value of
$S$ on a fixed positive form.  Thus $S$ is continuous, hence a current of
order $0$ relative to $(\mathscr D,K)$.
\end{proof}

\begin{remark}[The missing third case]
\label{rem:ch14-third-case}
The case not covered is that of a linear functional nonnegative only on
\emph{strongly} positive forms, which one would like to conclude is a weakly
positive current.  This is asserted in \cite[Proposition~7.2.9]{CLD25}; its
proof there dominates an arbitrary symmetric $\alpha$ by a multiple of
$\eta^p$ inside the strongly positive cone, using
\cite[Proposition~1.2.13]{CLD25}.  That proposition places $\eta^p$ in the
\emph{relative} interior of the strongly positive cone, so the resulting
estimate controls $S$ on the linear span of that cone; and by the warning
after \Cref{lem:ch14-eta-dominates} the span is proper as soon as $r\ge4$ and
$2\le p\le r-2$.  We therefore use only the two cases proved above.  They
suffice everywhere in this book: every current whose continuity we deduce
from positivity---measures, integration currents of cycles and skeleta,
fibre-integration functionals, curvature currents of psh
metrics---is nonnegative on positive forms, usually because it is strongly
positive.
\end{remark}

\begin{remark}
\Cref{prop:ch14-automatic-continuity} is the reason why, in this chapter, a
positive object is almost never checked for continuity: constructing a
positive functional constructs a current.  It is used below to see that
$\delta_X$, the integration current of a cycle, and the fibre-integration
functionals are currents, without any estimate.  The analogous statement in
the complex setting is classical; what is not classical is the proof, which
here rests on \Cref{lem:ch14-eta-dominates} and hence on the fact that a
tropical chart has a Euclidean target.
\end{remark}

\begin{proposition}
\label{prop:ch14-positive-and-negative}
Let $T$ be a current of bidimension $(p,p)$ on $X$, locally holomorphically
separated and paracompact.
\begin{enumerate}
\item If both $T$ and $-T$ are positive, then $T=0$.
\item There exist $X$, $p$ and a nonzero $T$ such that both $T$ and $-T$ are
  weakly positive.
\end{enumerate}
\end{proposition}

\begin{proof}
(1)  Let $\omega\in A_{X,c}^{p,p}$.  Since $T$ is symmetric, $\langle
T,\omega\rangle=\langle T,J\omega\rangle$, so $T$ annihilates every
antisymmetric form, and we may assume $\omega$ symmetric.  We claim that
$\omega$ is a difference of two positive compactly supported smooth forms.
Cover $\Supp\omega$ by finitely many chart domains, choose a subordinate
partition of unity $(\rho_i)$, and apply
\Cref{lem:ch14-eta-dominates} on each chart: with
$\Omega_i=\rho_if_{i,\trop}^*\eta_i^{p}$, which is strongly positive and
compactly supported, and with $c_i$ large enough,
$\rho_i\omega+c_i\Omega_i$ is positive.  Summing over $i$,
\[
 \omega=\Bigl(\omega+\sum_ic_i\Omega_i\Bigr)-\sum_ic_i\Omega_i
\]
exhibits $\omega$ as a difference of two positive compactly supported forms,
the subtrahend being even strongly positive.  Applying $T$ and $-T$, each term pairs
nonnegatively with $T$ and with $-T$, so $\langle T,\omega\rangle=0$.

(2)  Take $X=(\mathbf G_{m,K}^4)^{\An}$ with $K$ nontrivially valued,
$p=2$, and $T=[\gamma]$ where $\gamma$ is the weakly null superform of
\Cref{ex:ch14-weakly-null} below, pulled back along the tropicalization map.
By that example, $\gamma$ pairs to zero with every strongly positive
superform of complementary bidegree, so both $[\gamma]$ and $-[\gamma]$ are
weakly positive; and $\gamma\neq0$, so $[\gamma]\neq0$ by
\Cref{rem:ch14-form-injectivity}.
\end{proof}

\begin{warning}{Two hierarchies, one word}
\Cref{prop:ch14-positive-and-negative} says that the positive cone of
currents is \emph{salient}---it contains no line---while the weakly positive
cone is not.  The failure of saliency is precisely the phenomenon that
Berndtsson calls a \emph{weakly null direction}
(\Cref{sec:ch14-convex-bt}), and it is the reason why his convex
Bedford--Taylor products require an extra hypothesis in the weakly positive
case.  It is also the reason why ``$T\ge0$ and $T\le0$ imply $T=0$'' is a
theorem here and not a triviality.
\end{warning}

\begin{example}[Convex corners]
\label{ex:ch14-convex-corner-restated}
On the tropical line let $u(t)=\max\{0,t\}$.  Distributionally,
\[
                     \dprimeddprime[u]=\delta_0.
\]
Indeed, for a test function $\varphi$,
\[
 \langle\dprimeddprime[u],\varphi\rangle
 =\int_{\mathbf R}\max(0,t)\varphi''(t)\,\rd t=\varphi(0).
\]
The two integrations by parts have no sign change in the current convention.
This is the local model for curvature concentrated on a skeleton.
\end{example}

The computation deserves to be written once with every step visible, because
the same two integrations by parts recur in
\Cref{chap:cld-poincare-lelong,chap:cld-bedford-taylor}.  Here $n=1$, the
current $[u]$ has bidegree $(0,0)$, and
$\dprimeddprime[u]$ has bidegree $(1,1)$, so it pairs with test functions.
By \Cref{prop:current-differential-calculus},
\[
 \langle\dprimeddprime[u],\varphi\rangle
 =\langle[u],\dprimeddprime\varphi\rangle
 =\int_{\mathbf R}u(t)\,\varphi''(t)\,\rd t,
\]
using \eqref{eq:appJ-real-hessian} to write
$\dprimeddprime\varphi=\varphi''\,\dprime t\wedge\ddprime t$ and
\Cref{def:ch11-lattice-measure} to evaluate the integral of a top-degree
superform against the standard lattice on $\mathbf R$.  Two integrations by
parts, legitimate because $\varphi$ has compact support and $u$ is Lipschitz
with $u'=\one_{[0,\infty)}$ almost everywhere, give
\[
 \int_{\mathbf R}u\varphi''
 =-\int_{\mathbf R}u'\varphi'
 =-\int_0^{\infty}\varphi'
 =\varphi(0).
\]
No sign is created: the first parts formula carries $-1$, and the second is
the fundamental theorem of calculus, which carries the second $-1$.  The
answer is a positive measure, in accordance with the convexity of $u$ and
with \Cref{thm:convex-ddc-positive}.

\section{Worked examples on the Berkovich line}
\label{sec:ch14-line-examples}

Before continuing with the general theory we compute in the smallest
nontrivial case.  Throughout this section $k$ is algebraically closed and
\emph{trivially} valued, and
\[
 X=\mathbf P^{1,\An}_k .
\]
By \Cref{chap:berkovich} and \Cref{chap:skeleta}, $X$ is the compact space
obtained from the disjoint union
\[
 \coprod_{b\in B}\bigl(\{b\}\times[-\infty,0]\bigr),
 \qquad B=\mathbf P^1(k),
\]
by identifying all the points $(b,0)$ to a single point $\eta$, the trivial
absolute value on $k(\mathbf P^1)$, and topologized so that a set $V$ is open
if and only if each slice $V_b$ is open in $[-\infty,0]$ and, when
$\eta\in V$, all but finitely many branches lie entirely in $V$.  We write
$X_b$ for the branch of index $b$, $x_b=(b,-\infty)$ for its endpoint, which
is the rigid point $b$, and $t$ for the coordinate on a branch, normalized so
that $t(\eta)=0$, $t(x_b)=-\infty$, and $t=\log|T-b|$ in a coordinate $T$
vanishing at $b$.  Note the sign: $t$ is \emph{minus} the tropical coordinate
$-\log|T-b|$ of \Cref{chap:cld-forms}, chosen so that a branch reads
$[-\infty,0]$ with the Gauss point at the right-hand end.  Since the metric on
a branch is $|\rd t|$ either way, no Laplacian sign is affected; a reader who
prefers the book's default may replace $t$ by $-t$ throughout this section and
read the branches as $[0,+\infty]$.
This is a tropical space of level $1$, boundaryless, and its calibration is
the constant lattice $\mathbf Z$ on each branch.

We record the two structural facts we shall use, both immediate from the
definitions of \Cref{chap:tropical-spaces} transported by
\Cref{thm:analytic-is-tropical-space}.

\begin{lemma}
\label{lem:ch14-star-smooth}
\begin{enumerate}
\item A function $f$ on $X$ is smooth if and only if each $f_b:=f|_{X_b}$ is
  smooth on $[-\infty,0]$ and constant near $-\infty$, all but finitely many
  $f_b$ are constant on the whole branch, and the outgoing slopes satisfy the
  harmonicity relation $\sum_bf_b'(0)=0$.
\item If $f:U\to\mathbf R^m$ is a tropical chart with $\eta\in U$, then there
  is a finite $B_1\subset B$ such that $f$ is constant on $X_b\cap U$ for
  every $b\notin B_1$.
\item Every compact subset of $X\setminus\{\eta\}$ meets only finitely many
  branches.
\end{enumerate}
\end{lemma}

\begin{proof}
(1) and (2) are the definitions of a smooth function and of a paralinear map
on the star, read off from the description of its domains: a domain
containing $\eta$ contains all but finitely many complete branches, and a
paralinear function on such a domain is required to be constant on all but
finitely many of them.  The harmonicity relation at $\eta$ is the balancing
condition of \Cref{def:ch10-weighted-balanced} for the fan with one ray per
branch, all weights $1$.

(3) Let $K\subset X\setminus\{\eta\}$ be compact.  Its complement is open and
contains $\eta$, hence contains all but finitely many complete branches; so
$K$ meets only finitely many branches.
\end{proof}

\begin{example}[Dirac masses and the Gauss point]
\label{ex:ch14-line-diracs}
For any $x\in X$ the evaluation $f\mapsto f(x)$ is a positive linear
functional on $A_{X,c}^{0,0}$, hence a positive current $\delta_x$ of
bidimension $(0,0)$ by
\Cref{prop:ch14-automatic-continuity}\,(1); it is the Radon measure
$\delta_x$ under the identification of
\Cref{ex:ch14-forms-and-measures}.  Two features distinguish the endpoints
from $\eta$.

First, $\delta_{x_b}$ is a weak limit of currents defined by smooth forms.
Indeed, let $\omega_N$ be the smooth $(1,1)$-form supported in $X_b$ whose
density with respect to $\dprime t\wedge\ddprime t$ is a fixed bump of mass
$1$ concentrated near $t=-N$.  For any smooth $f$, part (1) of
\Cref{lem:ch14-star-smooth} makes $f_b$ constant on
$[-\infty,-N_f]$ for some $N_f$, so
$\langle[\omega_N],f\rangle=f_b(-\infty)=f(x_b)$ for $N\ge N_f$; the sequence
is eventually constant.

Second, $\delta_{x_b}$ is not $\dprimeddprime$ of anything smooth, and it is
not the restriction of a current of bidimension $(1,1)$: by
\Cref{cor:ch14-rigid-invisible} the rigid points support no current of
bidimension $(1,1)$ at all.  There is no contradiction: $\delta_{x_b}$ has
bidimension $(0,0)$.
\end{example}

\begin{example}[The Laplacian of a model function]
\label{ex:ch14-line-model-function}
Let $b_1,\ldots,b_r\in B$ be distinct and let $c_1,\ldots,c_r>0$.  Define
$u:X\to\mathbf R$ by $u=0$ outside $X_{b_1}\cup\cdots\cup X_{b_r}$ and, on
$X_{b_i}$,
\[
 u_{b_i}(t)=\max\{-c_i,\,c_it\}
 \qquad\text{(so }u(\eta)=0,\ u(x_{b_i})=-c_i).
\]
Then $u$ is continuous, is a difference of tropically convex functions on
each branch, and
\[
 \dprimeddprime[u]
 =\sum_{i=1}^rc_i\,\delta_{p_i}
   -\Bigl(\sum_{i=1}^rc_i\Bigr)\delta_\eta,
\]
where $p_i\in X_{b_i}$ is the point with $t=-1$.  Indeed, on each branch the
computation is that of \Cref{ex:ch14-convex-corner-restated}: a convex corner
of slope jump $c_i$ at $t=-1$ contributes $c_i\delta_{p_i}$, and at $\eta$
the outgoing slopes of $u$ are $c_i$ into the branch $b_i$ and $0$ elsewhere,
and \Cref{thm:ch11-graph-green}\,(1) together with
\Cref{lem:ch11-graph-trace} evaluates the atom at a vertex as the sum of the
outgoing slopes, here $\sum_i(-c_i)$.  The total mass is zero, as it must be for a
$\dprimeddprime$ of a function on a compact boundaryless space, by
\Cref{cor:analytic-stokes-green}.

This is the local model of the curvature of a model metric: compare
\Cref{ex:ch08-p1-tripod} and \eqref{eq:cld-curve-graph-laplacian}.
\end{example}

\begin{example}[A current with unbounded relative order]
\label{ex:ch14-star-unbounded-order}
Suppose $B$ is infinite and fix an injection
$\mathbf N_{\ge1}\hookrightarrow B$, $k\mapsto b_k$.  Define
\[
 \langle T,f\rangle=\sum_{k\ge1}f_{b_k}^{(k)}(-1),
 \qquad f\in A_{X,c}^{0,0}=C^\infty(X),
\]
where $f^{(k)}$ is the $k$-th derivative in the branch coordinate $t$.  Then:
\begin{enumerate}
\item the sum has only finitely many nonzero terms, so $T$ is a well-defined
  linear functional;
\item $T$ is a current of bidimension $(0,0)$;
\item $T$ has no global order: for every $m$ there is a chart datum relative
  to which $T$ is not of order $\le m$.
\end{enumerate}
\end{example}

\begin{proof}
(1) By \Cref{lem:ch14-star-smooth}\,(1), all but finitely many $f_b$ are
constant, and a constant function has vanishing derivatives of all positive
orders.

(2) Let $(\mathscr D,K)$ be a chart datum.  We claim there is a finite
$B_0\subset B$ such that every $\omega\in A^{0,0}_{\mathscr D,K}(X)$ is
constant on $X_b$ for $b\notin B_0$.  If $\eta\notin K$, take for $B_0$ the
finite set of branches meeting $K$, which is finite by
\Cref{lem:ch14-star-smooth}\,(3); a form supported in $K$ vanishes on the
other branches.  If $\eta\in K$, choose a chart $f_j:U_j\to\mathbf R^{m_j}$
of the finite family $\mathscr D$ with $\eta\in U_j$.  Let $B_1$ be the
finite exceptional set of \Cref{lem:ch14-star-smooth}\,(2) for $f_j$, and let
$B_2$ be the finite set of branches not entirely contained in $U_j$; both are
finite, the second because $U_j$ is an open neighbourhood of $\eta$.  Put
$B_0=B_1\cup B_2$.  For $b\notin B_0$ the branch $X_b$ lies in $U_j$ and
$f_j$ is constant on it, so any $\omega$ with
$\omega|_{U_j}=f_{j,\trop}^*\alpha_j$ is constant on $X_b$.

Set $m=\max\{k:b_k\in B_0\}$, a finite number, with $m=0$ if no $b_k$ lies in
$B_0$.  Then for $\omega\in A^{0,0}_{\mathscr D,K}(X)$,
\[
 |\langle T,\omega\rangle|
 =\Bigl|\sum_{k:\,b_k\in B_0}\omega_{b_k}^{(k)}(-1)\Bigr|
 \le\#B_0\cdot\max_{b\in B_0}\|\omega\|_{W_b,m},
\]
where $W_b\subset X_b$ is a compact neighbourhood of the point $t=-1$.  By
\Cref{prop:test-form-continuity-criterion}, $T$ is a current.

(3) Fix $m$ and let $N>m$.  Choose once and for all
$\chi\in C_c^\infty(\mathbf R)$ with $\Supp\chi\subset[-1,1]$ and
$\chi^{(N)}(0)=1$; for instance $\chi(u)=u^N\psi(u)/N!$ with $\psi$ a bump
equal to $1$ near $0$.  For $\lambda\ge1$ put
\[
 g_\lambda(t)=\lambda^{-N}\chi\bigl(\lambda(t+1)\bigr),
\]
regarded as a function on the branch $X_{b_N}$ and extended by $0$ to the
rest of $X$.  Each $g_\lambda$ is smooth on $X$: it is supported in
$\{|t+1|\le1/\lambda\}$, hence away from $\eta$ and away from $-\infty$, so
the constancy and harmonicity conditions of
\Cref{lem:ch14-star-smooth}\,(1) hold trivially.  All the $g_\lambda$,
$\lambda\ge1$, lie in the single chart datum $(\mathscr D,K)$ consisting of
the coordinate chart $t$ on the interior of $X_{b_N}$ together with
$K=\{|t+1|\le1\}$.

Now $g_\lambda^{(N)}(-1)=\lambda^{-N}\lambda^{N}\chi^{(N)}(0)=1$, so
$\langle T,g_\lambda\rangle=1$ for every $\lambda$, while for $j\le m<N$
\[
 g_\lambda^{(j)}(t)=\lambda^{j-N}\chi^{(j)}\bigl(\lambda(t+1)\bigr),
 \qquad\text{so}\qquad
 \|g_\lambda\|_{K,m}\le\lambda^{m-N}\|\chi\|_{C^m}
 \xrightarrow[\lambda\to\infty]{}0 .
\]
Hence no bound of the shape \eqref{eq:ch14-order-bound} with this $m$ can
hold on that datum.  As $m$ was arbitrary, $T$ has no global order.
\end{proof}

\begin{example}[An order-zero current that is not a measure]
\label{ex:ch14-star-not-measure}
Keep $B$ infinite and let $(c_b)_{b\in B}$ be real numbers.  Define
\[
 \langle S,f\rangle=\sum_{b\in B}c_b\bigl(f(\eta)-f(x_b)\bigr)
 =\sum_{b\in B}c_b\bigl(f_b(0)-f_b(-\infty)\bigr).
\]
Then $S$ is a current of bidimension $(0,0)$, of order $0$ relative to every
chart datum, and $S$ is the current attached to a signed Radon measure if and
only if $(c_b)$ is summable, in which case
\[
 S=\Bigl(\sum_bc_b\Bigr)\delta_\eta-\sum_bc_b\delta_{x_b}.
\]
In particular, for $(c_b)$ not summable, $S$ is an order-zero current that is
not a measure.
\end{example}

\begin{proof}
Well-definedness and continuity: by \Cref{lem:ch14-star-smooth}\,(1) all but
finitely many $f_b$ are constant, and for those
$f_b(0)-f_b(-\infty)=0$; the argument of
\Cref{ex:ch14-star-unbounded-order}\,(2) produces, for each chart datum, a
finite $B_0$ outside which every admissible $\omega$ is constant on the
branch, so
\[
 |\langle S,\omega\rangle|
 \le\Bigl(2\sum_{b\in B_0}|c_b|\Bigr)\|\omega\|_{C^0},
\]
an order-zero bound with a constant depending on the datum.

Suppose $(c_b)$ is summable.  Then the displayed measure is a finite signed
Radon measure, and testing it against a smooth $f$ reproduces
$\langle S,f\rangle$, all sums converging absolutely.

Suppose conversely that $S=[\mu]$ for a signed Radon measure $\mu$ on the
compact space $X$.  Then there is $C$ with $|\langle S,f\rangle|\le
C\|f\|_{C^0(X)}$ for every smooth $f$, since $|\mu|(X)<\infty$.  Fix a finite
$B_0\subset B$.  Choose $\varphi:[-\infty,0]\to[0,1]$ smooth with
$\varphi\equiv0$ on $[-1,0]$ and $\varphi\equiv1$ on $[-\infty,-2]$, and set
\[
 f_b=-\operatorname{sign}(c_b)\,\varphi\quad(b\in B_0),
 \qquad f_b=0\quad(b\notin B_0).
\]
This $f$ is smooth by \Cref{lem:ch14-star-smooth}\,(1): each $f_b$ is smooth
and constant near $-\infty$, all but finitely many are identically zero, and
every $f_b$ vanishes identically near $t=0$, so $f(\eta)=0$ and all the
slopes $f_b'(0)$ vanish, whence the harmonicity relation holds.  Moreover
$\|f\|_{C^0}\le1$ and $f(x_b)=-\operatorname{sign}(c_b)$ for $b\in B_0$.
Therefore
\[
 \langle S,f\rangle
 =\sum_{b\in B_0}c_b\bigl(0-(-\operatorname{sign}(c_b))\bigr)
 =\sum_{b\in B_0}|c_b|\le C .
\]
Letting $B_0$ increase gives $\sum_b|c_b|\le C<\infty$.
\end{proof}

\begin{warning}{Order zero does not mean ``measure coefficients''}
On a complex manifold, a current of order $0$ is a form whose coefficients
are Radon measures, by the Riesz representation theorem applied on each
compact set.  \Cref{ex:ch14-star-not-measure} shows that the corresponding
statement is \emph{false} on a Berkovich space, even a compact one: the
constant in the order-zero estimate depends on the chart datum, and no single
constant works for the compact space $X$.  The Riesz theorem is not being
contradicted; its hypothesis, boundedness for one supremum norm on
$C_c(X)$, simply fails.  This is a first, entirely elementary, instance of
the phenomena collected in \Cref{sec:ch14-compactness}.
\end{warning}

\begin{example}[Positive but not strongly positive]
\label{ex:ch14-positive-not-strong-current}
Let $K$ be nontrivially valued and let $X=(\mathbf G_{m,K}^4)^{\An}$, whose
tropicalization map $\trop:X\to\mathbf R^4$ is proper with image all of
$\mathbf R^4$, calibrated with multiplicity $1$, by
\Cref{prop:trop-torus-skeleton}.  Let $\omega\in A^{2,2}(\mathbf R^4)$ be the
constant-coefficient superform
\[
 \omega=-\bigl(\dprime x_{12}+\dprime x_{34}\bigr)
        \wedge\bigl(\ddprime x_{12}+\ddprime x_{34}\bigr)
\]
of \Cref{ex:ch09-positive-not-strong}, and set
$T=[\trop^*\omega]$, a current of bidegree $(2,2)$ and bidimension $(2,2)$
on $X$.  Then $T$ is positive but not strongly positive.
\end{example}

\begin{proof}
$T$ is positive by \Cref{prop:ch14-form-current-cones}, since $\omega$ is
positive by \Cref{ex:ch09-positive-not-strong} and pullback preserves the
positive cone (\Cref{prop:ch09-positivity-inclusions}).

Let $\beta\in A^{2,2}(\mathbf R^4)$ be the weakly positive superform of
\Cref{ex:ch09-weak-not-positive}, for which
$\omega\wedge\beta=-2\,\tau_4$.  Let $\chi\ge0$ be a smooth compactly
supported function on $\mathbf R^4$ with $\int_{\mathbf R^4}\chi\,\rd x=1$,
and put $\alpha=\trop^*(\chi\beta)$.  This is a compactly supported
weakly positive test form of bidegree $(2,2)$: compact support because
$\trop$ is proper, weak positivity because pullback preserves
$\mathrm{WP}_2$ and multiplication by $\chi\ge0$ preserves it.  Now
\[
 \langle T,\alpha\rangle
 =\int_X\trop^*\bigl(\chi\,\omega\wedge\beta\bigr)
 =-2\int_{\mathbf R^4}\chi\,\rd x=-2<0,
\]
the middle equality by \Cref{thm:cld-analytic-integration} and the
multiplicity-one calibration.  A strongly positive current is nonnegative on
weakly positive test forms; hence $T$ is not strongly positive.
\end{proof}

\begin{remark}
The example uses dimension $4$ and bidegree $(2,2)$ because, by the
discussion after \Cref{prop:ch09-positivity-inclusions}, the three cones
coincide for $p\in\{0,1,n-1,n\}$, hence for every $p$ when $n\le3$.  On
curves and surfaces the distinction between the three positivity notions for
currents is therefore invisible, which is exactly why it is so often
mislaid.
\end{remark}

\section{Two real-geometric incarnations of supercurrents}
\label{sec:ch14-berndtsson}

The algebra of superforms from \cref{chap:superforms} is not confined to tropical
geometry.  Two constructions of Berndtsson clarify both its reach and the
limits of the analogies used in this book.

\subsection{A metric-dependent current attached to a submanifold}
\label{sec:ch14-berndtsson-submanifold}

Let $E=\mathbf R^N$ carry a Euclidean scalar product, and write
$\beta=\sum_i \rd x_i\wedge \rd\xi_i$ in an orthonormal coordinate system.  If
$M\subset E$ is a smooth $m$-manifold of codimension $c$, choose locally an
orthonormal conormal frame $\nu_1,\ldots,\nu_c$ and let $\nu_i^\#$ be its copy
in the second exterior factor.  Berndtsson's supercurrent is, up to the
positive orientation constant fixed by the Berezin convention,
\[
 [M]_{\mathrm s}
 =c_c\,\nu_1\wedge\cdots\wedge\nu_c
       \wedge\nu_1^\#\wedge\cdots\wedge\nu_c^\#\,*\rd S_M .
\]
Here $*\rd S_M$ denotes the ordinary distributional integration density along
$M$ in the normal-current convention.  An orthogonal change of conormal
frame leaves the expression unchanged.
The current is symmetric and positive, and wedging a form with
$[M]_{\mathrm s}$ discards its normal components
\cite[\S4]{Ber19Super}.

In the notation of this book, $\nu_i^\#=J\nu_i$ and
$c_c=(-1)^{c(c-1)/2}$ is the constant \eqref{eq:appJ-cp}, so that
\[
 c_c\,\nu_1\wedge\cdots\wedge\nu_c\wedge J\nu_1\wedge\cdots\wedge J\nu_c
 =\nu_1\wedge J\nu_1\wedge\cdots\wedge\nu_c\wedge J\nu_c
\]
by \Cref{lem:ch09-generators}; this is the interlaced normal form, and it is
manifestly a strongly positive $(c,c)$-superform.  We first make the pairing
explicit; the computation is elementary and is what makes the comparison of
\Cref{prop:ch14-berndtsson-vs-tropical} possible.

\begin{lemma}
\label{lem:ch14-berndtsson-pairing}
Let $\alpha$ be a smooth compactly supported $(m,m)$-superform on $E$.  At a
point $x\in M$, let $e_1,\ldots,e_m$ be an orthonormal basis of $T_xM$ and
complete it by normals $e_{m+1},\ldots,e_N$; write $\theta_1,\ldots,\theta_N$
for the dual coframe, so that $\nu_i=\theta_{m+i}$.  Let $a(x)$ be defined by
\[
 i^*\alpha\bigr|_x
 =a(x)\,\dprime\theta_1\wedge\ddprime\theta_1\wedge\cdots
        \wedge\dprime\theta_m\wedge\ddprime\theta_m ,
\]
that is, $a$ is the density of the restriction $i^*\alpha$ with respect to
the Riemannian top orientation of $M$.  Then
\[
 \bigl\langle[M]_{\mathrm s},\alpha\bigr\rangle
 =\int_Ma\,\rd S_M .
\]
In particular $[M]_{\mathrm s}$ is the \emph{Riemannian} integration current
of $M$: it is symmetric, strongly positive, of order $0$, and its value on
$\alpha$ depends only on the tangential part of $\alpha$ along $M$.
\end{lemma}

\begin{proof}
Fix $x$ and expand
$\alpha=\sum_{|I|=|J|=m}\alpha_{IJ}\,\dprime\theta_I\wedge\ddprime\theta_J$
in the coframe.  The superform
$\prod_{i=1}^{c}\bigl(\nu_i\wedge J\nu_i\bigr)
=\prod_{i>m}\bigl(\dprime\theta_i\wedge\ddprime\theta_i\bigr)$
contains every index $>m$ in both blocks, so in the product with $\alpha$
only the term $I=J=\{1,\ldots,m\}$ survives.  Writing that term in interlaced
form by \Cref{lem:ch09-generators},
\[
 \alpha_{[m][m]}\,\dprime\theta_{[m]}\wedge\ddprime\theta_{[m]}
 =a\prod_{i\le m}\bigl(\dprime\theta_i\wedge\ddprime\theta_i\bigr),
 \qquad a=c_m\alpha_{[m][m]},
\]
we obtain, since each factor $\dprime\theta_i\wedge\ddprime\theta_i$ has even
total degree and therefore commutes with everything
(\Cref{prop:ch09-wedge-signs}),
\[
 \Bigl(\prod_{i>m}\dprime\theta_i\wedge\ddprime\theta_i\Bigr)\wedge\alpha
 =a\prod_{i=1}^{N}\bigl(\dprime\theta_i\wedge\ddprime\theta_i\bigr)
 =a\,\tau_N,
\]
with $\tau_N$ the positive top orientation
\eqref{eq:appJ-lagerberg-orientation}.  Integrating the density
$a$ against $\rd S_M$ gives the formula.  Symmetry follows because $J$ fixes
$\tau_N$ and permutes the two blocks of the normal factor; strong positivity
and order $0$ follow because the pairing is the integral of $a$, which is
nonnegative whenever $\alpha$ is weakly positive---by
\Cref{thm:ch09-positivity-duality} the tangential density of a weakly
positive form against a strongly positive one is nonnegative---and bounded
by $\|\alpha\|_{C^0}$ times the volume of $M\cap\Supp\alpha$.
\end{proof}

\begin{proposition}
\label{prop:ch14-berndtsson-metric-dependence}
Let $M\subset E$ be a smooth $m$-submanifold and let $g$ denote the ambient
scalar product used in the construction.
\begin{enumerate}
\item Replacing $g$ by $\lambda^2g$ for a constant $\lambda>0$ replaces
  $[M]_{\mathrm s}$ by $\lambda^{-m}[M]_{\mathrm s}$.
\item More generally, let $A\in\operatorname{GL}(N,\mathbf R)$ and let
  $g_A(u,v)=g(Au,Av)$.  Then, for $M$ a linear subspace,
  \[
   [M]_{\mathrm s}^{g_A}
   =\bigl|\!\det(A|_M)\bigr|^{-1}\,[M]_{\mathrm s}^{g},
  \]
  where $\det(A|_M)$ is computed with respect to $g$-orthonormal bases of $M$
  and of $AM$; in particular $[M]_{\mathrm s}$ is unchanged if and only if
  $A$ preserves the $m$-dimensional volume of $M$.
\end{enumerate}
\end{proposition}

\begin{proof}
By \Cref{lem:ch14-berndtsson-pairing}, $[M]_{\mathrm s}$ is determined by two
pieces of metric data: the Riemannian volume $\rd S_M$, and the Riemannian
top orientation of $M$ used to define the density $a$.  Under
$g\mapsto\lambda^2g$ the volume scales by $\lambda^{m}$, while a
$g$-orthonormal basis of $T_xM$ becomes $\lambda^{-1}$ times a
$\lambda^2g$-orthonormal one; hence the coframe scales by $\lambda$, the top
orientation $\prod_{i\le m}(\dprime\theta_i\wedge\ddprime\theta_i)$ scales by
$\lambda^{2m}$, and therefore $a$ scales by $\lambda^{-2m}$.  The product
$a\,\rd S_M$ scales by $\lambda^{m}\lambda^{-2m}=\lambda^{-m}$, which is (1).

For (2) the same bookkeeping applies with the scalar factor $\lambda$
replaced by the linear map $A|_M$: a $g_A$-orthonormal basis of $M$ is the
$A$-preimage of a $g$-orthonormal basis of $AM$, the $g_A$-volume of $M$ is
$|\det(A|_M)|$ times the $g$-volume, and the top orientation acquires
$|\det(A|_M)|^{2}$.  The two effects combine to
$|\det(A|_M)|^{1-2}=|\det(A|_M)|^{-1}$.
\end{proof}

\begin{warning}{$[M]_{\mathrm s}$ is not a metric-free invariant of $M$}
\Cref{prop:ch14-berndtsson-metric-dependence} shows that
$[M]_{\mathrm s}$ scales nontrivially under a homothety of the ambient
metric, whereas the calibrated tropical current $\delta_M$ of a rational
subspace, being defined by lattice measures, does not move at all.  Any
statement identifying the two therefore carries a hidden normalization; the
identification holds only in the unimodular situation described in
\Cref{prop:ch14-berndtsson-vs-tropical}.
\end{warning}

\subsubsection{Minimality}

This construction remembers the scalar product twice: it selects
orthonormal conormals and the Riemannian density $\rd S_M$.  Its derivative is
controlled by the second fundamental form.  If $\vec H$ is the mean-curvature
vector in Berndtsson's sign convention, his traced identity gives
\[
       \rd\bigl([M]_{\mathrm s}\wedge\beta^{m-1}\bigr)=0
       \quad\Longleftrightarrow\quad \vec H=0.
\]
Thus the traced supercurrent is closed exactly when $M$ is minimal
\cite[Corollary~5.2]{Ber19Super}.  The untraced $[M]_{\mathrm s}$ is generally
not closed; its differential retains the full second fundamental form.  The
paper also derives a variation formula for its trace under mean-curvature
flow.  We use the corrected 2019 version: the proof in Section~8 of the first
arXiv version was replaced, while the main results were unchanged.

We do not reprove the identification of $\rd([M]_{\mathrm s}\wedge\beta^{m-1})$
with the first-variation functional; that is the content of the citation just
given.  What we do prove is the classical statement to which it reduces, so
that the reader can see exactly which geometric fact is being asserted.
Recall that for $x\in M$ and an orthonormal tangent frame
$e_1,\ldots,e_m$ the second fundamental form is
$\mathrm{II}(u,v)=(\nabla_uv)^{\perp}$ and the \emph{mean-curvature vector}
is its unnormalized trace,
\[
 \vec H=\sum_{i=1}^m\mathrm{II}(e_i,e_i)
 =\sum_{i=1}^m\bigl(\nabla_{e_i}e_i\bigr)^{\perp},
\]
a normal field independent of the frame.  The manifold $M$ is \emph{minimal}
if $\vec H\equiv0$.

\begin{proposition}
\label{prop:ch14-first-variation}
Let $M\subset\mathbf R^N$ be a smooth properly embedded $m$-submanifold
without boundary and let $Y$ be a smooth compactly supported vector field on
$\mathbf R^N$.  Then
\[
 \int_M\operatorname{div}_MY\,\rd S_M
 =-\int_M\langle Y,\vec H\rangle\,\rd S_M ,
\]
where $\operatorname{div}_MY=\sum_{i=1}^m\langle\nabla_{e_i}Y,e_i\rangle$.
Consequently $\int_M\operatorname{div}_MY\,\rd S_M=0$ for every such $Y$ if
and only if $M$ is minimal.
\end{proposition}

\begin{proof}
Decompose $Y=Y^{\mathsf T}+Y^{\perp}$ into its tangential and normal parts
along $M$; both are smooth and compactly supported on $M$.  Then
$\operatorname{div}_MY=\operatorname{div}_MY^{\mathsf T}
+\operatorname{div}_MY^{\perp}$.

For the tangential part, $\operatorname{div}_MY^{\mathsf T}$ is the intrinsic
Riemannian divergence of a compactly supported vector field on the
boundaryless manifold $M$, so
$\int_M\operatorname{div}_MY^{\mathsf T}\rd S_M=0$ by the divergence theorem
on $M$.

For the normal part, differentiate the identity
$\langle Y^{\perp},e_i\rangle=0$, valid along $M$, in the direction $e_i$:
\[
 \langle\nabla_{e_i}Y^{\perp},e_i\rangle
 =-\langle Y^{\perp},\nabla_{e_i}e_i\rangle
 =-\bigl\langle Y^{\perp},(\nabla_{e_i}e_i)^{\perp}\bigr\rangle .
\]
Summing over $i$ gives
$\operatorname{div}_MY^{\perp}=-\langle Y^{\perp},\vec H\rangle
=-\langle Y,\vec H\rangle$, the last equality because $\vec H$ is normal.
Integrating proves the identity.

If $\vec H\equiv0$ the integral vanishes for all $Y$.  Conversely, if
$\vec H(x)\ne0$ for some $x$, choose $Y$ supported near $x$ with
$Y=\chi\vec H$ for a nonnegative bump $\chi$ equal to $1$ at $x$ (extending
$\vec H$ smoothly to a neighbourhood in $\mathbf R^N$); then
$-\int_M\chi|\vec H|^2\rd S_M<0$.
\end{proof}

\begin{remark}[What is imported and what is not]
\label{rem:ch14-berndtsson-import}
\Cref{prop:ch14-first-variation} is a statement about $M$ alone.  The
supercurrent enters through \cite[Corollary~5.2]{Ber19Super}, which
identifies $\langle \rd([M]_{\mathrm s}\wedge\beta^{m-1}),\cdot\rangle$, tested
against the forms determined by a compactly supported vector field $Y$, with
$\int_M\operatorname{div}_MY\,\rd S_M$ up to a positive constant.  Given
that identification, \Cref{prop:ch14-first-variation} is exactly the
equivalence displayed above.  We have separated the two halves because only
the second is specific to the superform formalism, and because the sign of
$\vec H$ is a convention that must be matched before the equivalence can be
used quantitatively; ours is the one for which the sphere of radius $R$ in
$\mathbf R^N$ has $\vec H$ pointing inward with $|\vec H|=(N-1)/R$.
\end{remark}

\subsubsection{Comparison with the calibrated tropical current}

This is not the calibrated tropical integration current.  A tropical cycle
uses integral tangent lattices, multiplicities, and balancing; no Euclidean
scalar product is part of its definition, and balancing supplies closedness.
Berndtsson's $[M]_{\mathrm s}$ instead applies to arbitrary smooth real
submanifolds, depends on the chosen metric, and measures failure of
minimality through its derivative.  The common language is the double
exterior algebra, not an identification of the two currents.

The relation, where both objects exist, can be written down exactly.  Recall
from \Cref{def:ch11-lattice-measure} that a rational subspace
$M\subset\mathbf R^N=N_{\mathbf R}$ with saturated tangent lattice
$N_M=N\cap M$ carries the lattice measure $\omega_M$, normalized so that a
lattice parallelotope has measure $1$, and that the associated integration
current is
\[
 \langle\delta_M,\alpha\rangle
 =\int_M\bigl(\text{density of }i^*\alpha
   \text{ w.r.t.\ the lattice top orientation}\bigr)\rd\omega_M .
\]

\begin{proposition}
\label{prop:ch14-berndtsson-vs-tropical}
Let $M\subset\mathbf R^N$ be a rational linear subspace of dimension $m$,
with saturated tangent lattice $N_M$, and let $g$ be a Euclidean scalar
product on $\mathbf R^N$.  Write
$\operatorname{covol}_g(N_M)$ for the $g$-volume of a fundamental
parallelotope of $N_M$ inside $M$.  Then
\[
 [M]_{\mathrm s}=\operatorname{covol}_g(N_M)^{-1}\,\delta_M .
\]
In particular the two currents agree if and only if $N_M$ is unimodular for
$g$, that is, has covolume $1$; and if $M$ is not rational, then $\delta_M$
is not defined at all while $[M]_{\mathrm s}$ is.
\end{proposition}

\begin{proof}
Fix a $g$-orthonormal basis $e_1,\ldots,e_m$ of $M$ and a lattice basis
$v_1,\ldots,v_m$ of $N_M$, and write $v_j=\sum_iA_{ij}e_i$, so that
$|\det A|=\operatorname{covol}_g(N_M)$.  Let $\theta_i$ and $\mu_j$ be the
dual coframes of $M^*$.  Then $\mu=(A^{-1})^{\mathsf T}\theta$, and since a
change of coframe with matrix $C$ multiplies a top-degree $(m,m)$-superform
by $(\det C)^2$---each of the two exterior blocks contributing one
factor---we get
\[
 \prod_{j\le m}\bigl(\dprime\mu_j\wedge\ddprime\mu_j\bigr)
 =(\det A)^{-2}\prod_{i\le m}
   \bigl(\dprime\theta_i\wedge\ddprime\theta_i\bigr).
\]
Consequently, if $i^*\alpha=a\prod_i(\dprime\theta_i\wedge\ddprime\theta_i)$
is the Riemannian expression of \Cref{lem:ch14-berndtsson-pairing}, then its
lattice expression is
$i^*\alpha=b\prod_j(\dprime\mu_j\wedge\ddprime\mu_j)$ with
$b=(\det A)^{2}a$.  On the other hand
$\omega_M=|\det A|^{-1}\rd S_M$, because the lattice parallelotope has
$g$-volume $|\det A|$ and lattice measure $1$.  Therefore
\[
 \begin{aligned}
 \langle\delta_M,\alpha\rangle
 &=\int_Mb\,\rd\omega_M
   =\int_M(\det A)^2a\,|\det A|^{-1}\rd S_M\\
 &=|\det A|\int_Ma\,\rd S_M
   =\operatorname{covol}_g(N_M)\,\langle[M]_{\mathrm s},\alpha\rangle,
 \end{aligned}
\]
using \Cref{lem:ch14-berndtsson-pairing} for the last equality.
\end{proof}

\begin{example}[The diagonal line: the two currents differ by $\sqrt2$]
\label{ex:ch14-diagonal-line}
Take $N=2$ with the standard lattice $\mathbf Z^2$ and standard scalar
product, and let $M=\mathbf R\cdot(1,1)$, a tropical cycle of dimension $1$
and weight $1$.  Its saturated tangent lattice is $\mathbf Z\cdot(1,1)$, of
covolume $\sqrt2$.  We verify $[M]_{\mathrm s}=\tfrac1{\sqrt2}\delta_M$
directly.

The unit conormal is $\nu=(\dprime x_1-\dprime x_2)/\sqrt2$, so
\[
 [M]_{\mathrm s}
 =\tfrac12\bigl(\dprime x_1-\dprime x_2\bigr)
   \wedge\bigl(\ddprime x_1-\ddprime x_2\bigr)\otimes\rd S_M .
\]
Let $\alpha=a\,\dprime x_1\wedge\ddprime x_1$.  In the expansion of
$\nu\wedge J\nu$, every term containing $\dprime x_1$ or $\ddprime x_1$ dies
against $\alpha$, leaving
\[
 \nu\wedge J\nu\wedge\alpha
 =\tfrac a2\,\dprime x_2\wedge\ddprime x_2\wedge
   \dprime x_1\wedge\ddprime x_1
 =\tfrac a2\,\tau_2 ,
\]
the reordering being free because the blocks have even total degree.
Parametrizing $M$ by $t\mapsto(t,t)$ we have $\rd S_M=\sqrt2\,\rd t$, so
\[
 \langle[M]_{\mathrm s},\alpha\rangle
 =\int_{\mathbf R}\tfrac{a(t,t)}2\sqrt2\,\rd t
 =\tfrac1{\sqrt2}\int_{\mathbf R}a(t,t)\,\rd t .
\]
On the other hand $i^*\alpha=a\,\dprime t\wedge\ddprime t$ and the lattice
measure is $\rd t$, because $t$ increases by $1$ across the lattice
generator $(1,1)$; hence
$\langle\delta_M,\alpha\rangle=\int a(t,t)\rd t$.  The ratio is
$1/\sqrt2=\operatorname{covol}(\mathbf Z(1,1))^{-1}$, as predicted.  The same
computation with $M=\mathbf R\cdot(1,0)$ gives ratio $1$: there the lattice
is unimodular and the two currents coincide.

Finally, let $M'=\mathbf R\cdot(1,\sqrt2)$.  Then $[M']_{\mathrm s}$ is a
perfectly good positive closed supercurrent of bidegree $(1,1)$ on
$\mathbf R^2$, while $M'$ carries no tropical cycle structure whatsoever: it
has no nonzero rational tangent vector, so $N\cap M'=0$ and
\Cref{def:ch11-lattice-measure} produces no measure.  No amount of
renormalization relates the two.
\end{example}

\begin{remark}[Why closedness is not the same phenomenon]
\label{rem:ch14-closedness-two-reasons}
Both $\delta_M$ and $[M]_{\mathrm s}$ are closed when $M$ is a linear
subspace, but for different reasons.  For $\delta_M$ closedness is
\Cref{thm:ch11-balancing-iff-stokes}: a single cell has no interior
codimension-one face, and a balanced complex has cancelling ones.  For
$[M]_{\mathrm s}$ it is \Cref{prop:ch14-first-variation}: a linear subspace
is totally geodesic, hence minimal.  For a curved $M$ the second reason
fails while the first is not even available, and for an irrational
polyhedral complex the first fails while the second may hold.  The two
notions of ``closed integration current'' are therefore genuinely different
theorems about the same word.
\end{remark}

\subsection{Convex Bedford--Taylor products and homogeneous currents}
\label{sec:ch14-convex-bt}

The convex theory of \cite{Ber26Convex} sharpens a point already visible in
\Cref{prop:current-positivity-calculus}.  In intermediate
bidegrees there are nonzero forms that are both weakly positive and weakly
negative; Berndtsson calls them \emph{weakly null}.  Consequently weak
positivity alone does not control every coefficient.  In that paper a form
is called \emph{strong} when it lies in the linear span of the strongly
positive cone.  This use of ``strong'' is linear-algebraic and should not be
confused with ``strongly positive.''

\subsubsection{Weakly null directions, explicitly}

\begin{definition}[Weakly null]
\label{def:ch14-weakly-null}
\index{weakly null form}\index{form!weakly null}
A symmetric superform $\gamma\in A^{p,p}(V)$ is \emph{weakly null} if both
$\gamma$ and $-\gamma$ are weakly positive, that is, if
$\gamma\wedge\beta=0$ for every $\beta\in\mathrm{SP}_{n-p}(V)$.  A symmetric
current $T$ of bidimension $(p,p)$ is weakly null if both $T$ and $-T$ are
weakly positive.
\end{definition}

By \Cref{prop:ch09-positivity-dictionary}\,(3) and
\Cref{thm:ch09-positivity-duality}, $\gamma$ is weakly null exactly when
$w^{\mathsf T}M(\gamma)w=0$ for every decomposable $w\in\Gamma_p$.  The
smallest instance is again $n=4$, $p=2$, and it is produced by the Pl\"ucker
quadratic form.

\begin{example}[A nonzero weakly null $(2,2)$-form on $\mathbf R^4$]
\label{ex:ch14-weakly-null}
Let $V=\mathbf R^4$, $p=2$, and let
\[
 P(w)=w_{12}w_{34}-w_{13}w_{24}+w_{14}w_{23}
\]
be the Pl\"ucker form of \Cref{ex:ch09-weak-not-positive}, with symmetric
matrix $\widehat P$, so that $w^{\mathsf T}\widehat Pw=P(w)$.  Define
$\gamma\in A^{2,2}(V)$ by $M(\gamma)=\widehat P$, that is,
\[
\begin{aligned}
 \gamma=-\tfrac12\Bigl[
  &\bigl(\dprime x_{12}\wedge\ddprime x_{34}
        +\dprime x_{34}\wedge\ddprime x_{12}\bigr)
  -\bigl(\dprime x_{13}\wedge\ddprime x_{24}
        +\dprime x_{24}\wedge\ddprime x_{13}\bigr)\\
  &+\bigl(\dprime x_{14}\wedge\ddprime x_{23}
        +\dprime x_{23}\wedge\ddprime x_{14}\bigr)\Bigr],
\end{aligned}
\]
the overall factor being $c_2=-1$.  Then $\gamma\neq0$ and $\gamma$ is weakly
null.  Moreover, in the notation of \Cref{ex:ch09-weak-not-positive}, the
weakly positive form there is $\beta=\delta_{\mathrm{Id}}-4\gamma$, where
$M(\delta_{\mathrm{Id}})=\mathrm{Id}$.
\end{example}

\begin{proof}
$\gamma\ne0$ because $\widehat P\ne0$ and $M$ is a linear isomorphism.  If
$w\in\Gamma_2$ is decomposable, say $w=\eta_1\wedge\eta_2$, then
$w\wedge w=0$; by the computation of \Cref{ex:ch09-weak-not-positive},
$w\wedge w=2P(w)\lambda$, so $P(w)=0$, that is,
$w^{\mathsf T}M(\gamma)w=0$.  By
\Cref{prop:ch09-positivity-dictionary}\,(3) both $\gamma$ and $-\gamma$ are
weakly positive.  The last assertion is the identity
$M(\beta)=\mathrm{Id}-4\widehat P$ recorded in
\Cref{ex:ch09-weak-not-positive}.
\end{proof}

\begin{corollary}
\label{cor:ch14-sp-not-spanning}
For $V=\mathbf R^4$ and $p=2$, the linear span of $\mathrm{SP}_2(V)$ is a
proper subspace of the space of symmetric $(2,2)$-forms; equivalently,
$\mathrm{SP}_2(V)$ has empty interior there.  The same holds for every
$n\ge4$ and $2\le p\le n-2$.
\end{corollary}

\begin{proof}
For $v\in\Gamma_2$ we have
$\langle M(\gamma),vv^{\mathsf T}\rangle=v^{\mathsf T}M(\gamma)v=0$, so the
nonzero symmetric matrix $M(\gamma)$ is orthogonal, for the trace pairing, to
every generator of $M(\mathrm{SP}_2(V))$ as described in
\Cref{prop:ch09-positivity-dictionary}\,(2), hence to its span.  For general
$n$ and $p$, embed $\mathbf R^4$ as the span of four of the coordinates and
wedge $\gamma$ with a strongly positive form in $p-2$ of the remaining
coordinates; the resulting form is again weakly null by
\Cref{prop:ch09-positivity-products}\,(3) applied to $\pm\gamma$.
\end{proof}

\begin{corollary}
\label{cor:ch14-weakly-null-currents}
Let $K$ be nontrivially valued and $X=(\mathbf G_{m,K}^4)^{\An}$.  The
current $[\trop^*\gamma]$ of bidegree $(2,2)$ is nonzero and weakly null.  By
contrast, by \Cref{prop:ch14-positive-and-negative}\,(1), no nonzero current
is both positive and negative.
\end{corollary}

\begin{proof}
Weak nullity follows from \Cref{prop:ch14-form-current-cones} applied to
$\pm\trop^*\gamma$, both of which are weakly positive by
\Cref{ex:ch14-weakly-null} and \Cref{prop:ch09-positivity-inclusions}.
Nonvanishing is \Cref{rem:ch14-form-injectivity}.
\end{proof}

\subsubsection{The convex Bedford--Taylor product}

For finite convex functions $\phi_i$ on a Euclidean domain, Berndtsson
constructs
\[
       \rd\rd^\#\phi_1\wedge\cdots\wedge \rd\rd^\#\phi_q\wedge\Omega
\]
when $\Omega$ is closed and either positive, or weakly positive and strong.
The product remains in the corresponding cone and is jointly continuous for
local uniform convergence of the convex potentials and weak convergence of
$\Omega$ \cite[Theorem~4.2]{Ber26Convex}.  The strong hypothesis in the
weakly positive case annihilates the weakly-null directions.  This is a
second reason not to assert that arbitrary positive-looking intermediate
forms may always be wedged while preserving positivity.

It is worth isolating the logical shape of that theorem, because the same
shape recurs in \Cref{chap:cld-bedford-taylor}.

\refstepcounter{theorem}\label{thm:ch14-convex-bt-imported}
\begin{quote}
\textbf{Foundational input \thetheorem\ (Berndtsson's convex
Bedford--Taylor product; source boundary).}
Let $U\subset\mathbf R^N$ be open and convex, let
$\phi_1,\ldots,\phi_q$ be finite convex functions on $U$, and let $\Omega$ be
a closed supercurrent of some bidegree $(a,a)$ with $a+q\le N$ on $U$, which
is either
positive, or weakly positive and strong in the sense above.  Then the product
\[
 \dprimeddprime\phi_1\wedge\cdots\wedge\dprimeddprime\phi_q\wedge\Omega
\]
is defined as the limit of the corresponding products for smooth convex
regularizations of the $\phi_i$; it is independent of the regularizations, it
lies in the same cone as $\Omega$, it is closed, and it is jointly continuous
for local uniform convergence of the $\phi_i$ and weak convergence of
$\Omega$.  This is \cite[Theorem~4.2]{Ber26Convex}, translated into the
present conventions by \Cref{rem:ch14-convention-firewall}.  This entire
statement, including existence, independence, positivity, closedness, and
joint continuity, is imported from the cited theorem; no part of its proof
is claimed here.  It is one of the three foundational inputs listed in
\Cref{sec:ch14-prerequisites}.
\end{quote}

What we can say without that input is why the hypothesis on $\Omega$ cannot
simply be ``weakly positive''.

\begin{proposition}
\label{prop:ch14-why-strong}
Let $\Omega$ be weakly null and nonzero, as in \Cref{ex:ch14-weakly-null}.
Then $\Omega$ is weakly positive, the wedge of $\Omega$ with any strongly
positive form of complementary bidegree is zero, and hence no
inequality of the shape
\[
 \bigl|\langle\Omega,\alpha\rangle\bigr|
 \le C\,\langle\Omega,\Theta\rangle\,\|\alpha\|_{C^0}
\]
with $\Theta$ strongly positive can control $\Omega$: the right-hand side
vanishes identically while the left-hand side does not.  Consequently the
mass estimates that drive the Bedford--Taylor construction---all of which
compare a product against a fixed strongly positive form---carry no
information for a weakly null $\Omega$, and the construction must exclude
those directions by hypothesis.
\end{proposition}

\begin{proof}
By \Cref{def:ch14-weakly-null} the pairing of $\Omega$ with any strongly
positive form of complementary bidegree is zero, so
$\langle\Omega,\Theta\rangle=0$ for every strongly positive $\Theta$.  On the
other hand $\Omega\ne0$, so by \Cref{thm:ch09-positivity-duality} there is a
symmetric $\alpha$ with $\langle\Omega,\alpha\rangle\ne0$; such an $\alpha$
necessarily fails to be strongly positive, and indeed fails to lie in the
span of the strongly positive cone by \Cref{cor:ch14-sp-not-spanning}.
\end{proof}

\subsubsection{Strongly homogeneous currents}

For $q=N$ and $\Omega=1$, the measure
$(\rd\rd^\#\phi)^N/N!$ agrees, with the Euclidean volume normalization, with the
Alexandrov Monge--Amp\`ere measure defined by the subgradient image.  If
$h_K$ is the support function of a convex body, the currents $\rd\rd^\#h_K$ are
strongly homogeneous: they are closed, annihilated by the appropriate
radial contraction, and invariant under positive scaling.  Mixed products
of such currents encode mixed volumes.  The resulting bilinear form has the
Lorentzian signature behind the Alexandrov--Fenchel inequality
\cite[\S\S5--6 and 10]{Ber26Convex}.  This is the real-convex counterpart of
the toric dictionary in \Cref{sec:toric-convex}; \Cref{chap:cld-bedford-taylor}
gives the local CLD product, while the convexity appendix returns to support
functions and mixed volumes.

The appendix treats support functions, mixed volumes, and the identity
\eqref{eq:appD-current-mixed-volume} in detail; see
\Cref{sec:appD-supercurrents} and
\Cref{thm:appD-support-current-mixed-volume}.  Here we prove only the two
current-theoretic statements that the appendix records as
\eqref{eq:appD-strong-homogeneity}, because they are statements about
currents and their proofs belong in this chapter.

We first extend the contraction operators of \Cref{def:ch09-contraction} to
currents.  Let $E$ denote the Euler, or radial, vector field on
$V=\mathbf R^N$, so $E(x)=x$; for a superform $\alpha$ put
$(\iota'_E\alpha)(x)=\iota'_{x}\bigl(\alpha(x)\bigr)$, and similarly for
$\iota''_E$.  These are $\mathbf R$-linear operators
$A^{p,q}\to A^{p-1,q}$, respectively $A^{p,q}\to A^{p,q-1}$, continuous for
each $C^m$ seminorm on a compact set, since their coefficients are
polynomials of degree $1$ in $x$.

\begin{definition}[Radial contraction of a current]
\label{def:ch14-radial-contraction}
For $T\in D^{p,q}(V)$ define $\iota'_ET\in D^{p+1,q}(V)$ and
$\iota''_ET\in D^{p,q+1}(V)$ by
\[
 \langle\iota'_ET,\alpha\rangle=\langle T,\iota'_E\alpha\rangle,
 \qquad
 \langle\iota''_ET,\alpha\rangle=\langle T,\iota''_E\alpha\rangle .
\]
\end{definition}

Pullback of currents along an \emph{invertible} affine map is available, and
we record its normalization once, since it is the only pullback of currents
used in this book; see \Cref{sec:ch14-pullback-failure} for why no more can
be expected.

\begin{definition}[Pullback along an affine isomorphism]
\label{def:ch14-affine-pullback}
Let $F:V\to V$ be an affine isomorphism with linear part $A$.  For
$T\in D^{p,q}(V)$ set
\[
 \langle F^*T,\alpha\rangle
 =|\det A|\,\bigl\langle T,(F^{-1})^*\alpha\bigr\rangle,
 \qquad\alpha\in A_c^{N-p,N-q}(V).
\]
\end{definition}

\begin{lemma}
\label{lem:ch14-affine-pullback-compat}
With this normalization, $F^*[u]=[F^*u]=[u\circ F]$ for every locally
integrable function $u$, and more generally $F^*[\omega]=[F^*\omega]$ for
every locally integrable form $\omega$.  Moreover $F^*$ commutes with
$\dprime$, $\ddprime$ and $J$, and $(F\circ G)^*=G^*F^*$.
\end{lemma}

\begin{proof}
Write $\alpha=a\,\tau_N$ for a top-degree test form.  Since
$(F^{-1})^*\tau_N=(\det A)^{-2}\tau_N$---each of the $2N$ covector factors
contributing one power of $A^{-1}$, and the two blocks contributing
$\det(A^{-1})$ each---we get
\[
 \bigl\langle[u],(F^{-1})^*\alpha\bigr\rangle
 =(\det A)^{-2}\int_Vu(y)\,a(A^{-1}y)\,\rd y ,
\]
while the affine change of variables $y=F(x)$ gives
\[
 \bigl\langle[u\circ F],\alpha\bigr\rangle
 =\int_Vu(F(x))a(x)\rd x
 =|\det A|^{-1}\int_Vu(y)a(A^{-1}y)\rd y .
\]
Comparing, we obtain
\[
 \begin{aligned}
  \langle[u\circ F],\alpha\rangle
   &=|\det A|^{-1}(\det A)^{2}
     \langle[u],(F^{-1})^*\alpha\rangle\\
   &=|\det A|\,\langle[u],(F^{-1})^*\alpha\rangle,
 \end{aligned}
\]
which is the assertion.  The
case of a general form is identical with $a$ replaced by the coefficient of
$\omega\wedge\alpha$.  Compatibility with the operators follows because
$(F^{-1})^*$ commutes with them on forms
(\Cref{prop:ch09-pullback-functorial}) and the constant $|\det A|$ does not
depend on the bidegree; functoriality is the multiplicativity of
$|\det|$.
\end{proof}

For $t>0$ we write $s_t(x)=tx$, so that $s_t^*$ is defined by
\Cref{def:ch14-affine-pullback} with $|\det A|=t^{N}$.

\begin{definition}[Strongly homogeneous]
\label{def:ch14-strongly-homogeneous}
\index{current!strongly homogeneous}
A current $T\in D^{p,p}(V)$ is \emph{strongly homogeneous} if it is closed,
$s_t^*T=tT$ for every $t>0$, and $\iota'_ET=\iota''_ET=0$.
\end{definition}

\begin{proposition}
\label{prop:ch14-strongly-homogeneous}
Let $h:V\to\mathbf R$ be convex and positively homogeneous of degree $1$, for
instance the support function $h_K$ of a nonempty compact convex set $K$.
Then $T_h=\dprimeddprime[h]$ is a closed positive current of bidegree
$(1,1)$, and it is strongly homogeneous.
\end{proposition}

\begin{proof}
$h$ is convex and finite, hence locally Lipschitz and locally integrable, so
$[h]$ is a current of bidegree $(0,0)$ and $T_h=\dprimeddprime[h]$ is defined.
It is positive by \Cref{thm:ch09-distributional-convexity} and closed because
$\dprime T_h=\dprime\dprimeddprime[h]=0$ by
\Cref{prop:current-differential-calculus}, and likewise for $\ddprime$.

\emph{Scaling.}  Homogeneity of degree $1$ says $h\circ s_t=t\,h$, so
\Cref{lem:ch14-affine-pullback-compat} gives
$s_t^*[h]=[h\circ s_t]=t[h]$.  Since $s_t$ is affine, $s_t^*$ commutes with
$\dprime$ and $\ddprime$ on currents by the same lemma, whence
\[
 s_t^*T_h=s_t^*\dprimeddprime[h]=\dprimeddprime s_t^*[h]
        =\dprimeddprime\bigl(t[h]\bigr)=t\,T_h .
\]

\emph{Contraction.}  We compute on smooth $h$ first and then pass to the
limit.  Suppose $h$ is smooth away from the origin and homogeneous of degree
$1$.  Then $\dprimeddprime h=\sum_{i,j}h_{ij}\dprime x_i\wedge\ddprime x_j$
by \eqref{eq:appJ-real-hessian}, and by
\Cref{def:ch09-contraction},
\[
 \iota'_E\bigl(\dprimeddprime h\bigr)
 =\sum_{i,j}h_{ij}(x)\,x_i\,\ddprime x_j .
\]
Differentiating the homogeneity relation
$\partial_jh(tx)=\partial_jh(x)$, valid for all $t>0$ because $\nabla h$ is
homogeneous of degree $0$, with respect to $t$ at $t=1$ gives
$\sum_ix_ih_{ij}(x)=0$ for every $j$.  Hence
$\iota'_E(\dprimeddprime h)=0$, and $\iota''_E(\dprimeddprime h)=0$ follows by
applying $J$ and using $J\iota'_EJ=\iota''_E$
(\Cref{lem:ch09-contraction-basic}\,(3)) together with the symmetry
$J\dprimeddprime h=\dprimeddprime h$.

For general $h$, Euler's identity holds distributionally.  Indeed, $h$ is
locally Lipschitz, the identity $\sum_i x_i\partial_i h=h$ holds almost
everywhere by positive homogeneity, and hence it holds as an identity of
distributions.  Applying $\partial_j$ distributionally gives
\[
 \sum_i x_i\partial_i\partial_j[h]=0
 \qquad (j=1,\ldots,N).
\]
Multiplication by $x_i$ and distributional differentiation are continuous,
so this calculation needs no homogeneous smoothing.  In the coefficient
description of $\dprimeddprime[h]$, the displayed identities say precisely
that $\iota'_ET_h=0$.  Finally $J T_h=T_h$, and
$J\iota'_EJ=\iota''_E$, so $\iota''_ET_h=0$ as well.
\end{proof}

\begin{remark}
\label{rem:ch14-homogeneous-appendix}
\Cref{prop:ch14-strongly-homogeneous} is the current-theoretic content of
\eqref{eq:appD-support-current} and
\eqref{eq:appD-strong-homogeneity}.  The mixed-volume identity
\eqref{eq:appD-current-mixed-volume},
\[
 \tfrac1{N!}\,T_{K_1}\wedge\cdots\wedge T_{K_N}
 =V(K_1,\ldots,K_N)\,\delta_0,
\]
is proved in \Cref{thm:appD-support-current-mixed-volume}; its only analytic
input is \Cref{thm:ch14-convex-bt-imported}, and its combinatorial input is
\Cref{lem:appD-support-calculus}.  The Lorentzian signature of the resulting
bilinear form is the Alexandrov--Fenchel inequality
\Cref{thm:appD-af}.  We do not repeat any of this here; the point of the
present subsection is that the word ``strongly homogeneous'' refers to three
separate conditions---closedness, degree-one scaling, and annihilation by
both radial contractions---and that all three are theorems, not
normalizations.
\end{remark}

\begin{warning}{Homogeneity is not a licence to multiply}
The phrase ``strongly homogeneous'' does not license products of arbitrary
weakly positive currents, and
\Cref{prop:ch14-strongly-homogeneous} produces only currents of bidegree
$(1,1)$, where by \Cref{prop:ch09-positivity-inclusions} the three positivity
notions coincide anyway.  The products
$T_{K_1}\wedge\cdots\wedge T_{K_q}$ of intermediate bidegree exist because
each factor is a $\dprimeddprime$ of a convex function, so
\Cref{thm:ch14-convex-bt-imported} applies; they do not exist for a general
strongly homogeneous pair.
\end{warning}

\begin{remark}[Convention firewall]
\label{rem:ch14-convention-firewall}
The formulas in \cite{Ber26Convex} cannot be copied symbol for symbol into
the Lagerberg--CLD convention.  In its Section~2,
\[
       J(\rd x_i)=-\rd\xi_i,\qquad J(\rd\xi_i)=\rd x_i,
       \qquad \rd^\#=J^{-1}\rd J,
\]
and top-degree integration uses a specified Berezin sign.  On functions,
$\rd\rd^\#\phi$ has Hessian coefficients in the ordered basis
$\rd x_i\wedge \rd\xi_j$, so it corresponds to our
$\dprimeddprime\phi$ after identifying $\rd x_i$ with $\dprime x_i$ and
$\rd\xi_i$ with $\ddprime x_i$.  Identities involving $J$, symmetry, current bidegrees, or
integration signs must nevertheless be recomputed.  Moreover, the 2019 and
2026 Berndtsson papers themselves use opposite signs for $J$.  Throughout
this book, the conventions of \Cref{chap:cld-forms} remain in force.
\end{remark}

\section{Restriction, support, and extension}

Currents form a sheaf on locally differentially separated tropical spaces,
hence on the analytic spaces considered here.  The support $\Supp(T)$ is the
complement of the largest open set on which $T$ vanishes.

\begin{proposition}
\label{prop:form-restriction-dimension-gap}
Let $i:Z\hookrightarrow X$ be an analytic subspace.  If the local tropical
dimension of $Z$ is everywhere smaller than $\max\{r,s\}$, then
$i^*\alpha=0$ for every $\alpha\in A_X^{r,s}$.
\end{proposition}
\begin{proof}
After passing to a tropical chart, the pullback is the restriction of a
superform to a polyhedron of dimension smaller than $r$ or smaller than
$s$.  Its alternating $r$- or $s$-covector factor therefore restricts to
zero.
\end{proof}

This restriction statement does not say that $\alpha$ vanishes on a
neighborhood of $Z$.  Nor does it imply that a current on $X\setminus Z$
extends, or that an extension is unique: even when an extension exists, one
may sometimes add a nonzero current supported on $Z$.  Removability requires
additional hypotheses, such as the precise growth, positivity, or canonical
extension results used in
\Cref{chap:cld-poincare-lelong,chap:cld-bedford-taylor}.

\subsection{Support theorems}
\label{sec:ch14-support-theorems}

We collect the statements that let one recognize a current from its support.
The first is the exact analogue of the classical fact that a distribution
supported at a point is a finite combination of derivatives of a Dirac mass;
the non-Archimedean version is stronger, because \emph{no} derivative of a
Dirac mass survives in high bidimension.

\begin{theorem}
\label{thm:ch14-support-dimension}
Let $X$ be pure $n$-dimensional and let $Z\subset X$ be closed with
$d_{\trop}(x)<\max\{r,s\}$ for every $x\in Z$.  Then no nonzero current of
bidimension $(r,s)$ has support contained in $Z$.  Equivalently, restriction
$D_{r,s}(X)\to D_{r,s}(X\setminus Z)$ is injective, and by
\Cref{lem:ch14-extension-by-zero} it is even bijective.
\end{theorem}

\begin{proof}
Let $\Supp T\subset Z$ and let $\alpha\in A_{X,c}^{r,s}$.  By
\Cref{lem:ch14-extension-by-zero} the form $\alpha$ has support in
$X\setminus Z$, hence disjoint from $\Supp T$, so
$\langle T,\alpha\rangle=0$ by \Cref{def:ch14-support}.
\end{proof}

\begin{theorem}
\label{thm:ch14-support-pushforward}
Let $g:Y\to X$ be a morphism of topologically separated analytic spaces and
let $T\in D_{r,s}(Y)$ have $g$-compact support, so that $g_*T$ is defined as
in \Cref{sec:ch14-pushforward}.  Then
\[
 \Supp(g_*T)\subset \overline{g\bigl(\Supp T\bigr)}.
\]
If $g$ is topologically proper, $g(\Supp T)$ is closed, so the closure may
be omitted.
\end{theorem}

\begin{proof}
Let $x\notin g(\Supp T)$.  If $g$ is proper, $g(\Supp T)$ is closed, so there
is an open $U\ni x$ with $U\cap g(\Supp T)=\emptyset$; then for
$\alpha$ compactly supported in $U$ the form $g^*\alpha$ is supported in
$g^{-1}(U)$, which misses $\Supp T$, so
$\langle g_*T,\alpha\rangle=\langle T,g^*\alpha\rangle=0$.  Hence
$g_*T$ vanishes on $U$.  In the merely $g$-compact case, replace
$g(\Supp T)$ by its closure and run the same argument.
\end{proof}

\begin{proposition}
\label{prop:ch14-support-fibre}
Let $g:Y\to X$ be a morphism and let $T\in D_{r,s}(Y)$ be supported in a
single fibre $Y_x$.  If $\dim Y_x<\max\{r,s\}$, then $T=0$.
\end{proposition}

\begin{proof}
A fibre is a weakly analytic subset of $Y$ in the sense of
\Cref{def:trop-weakly-analytic}, so its tropical dimension at each point is
at most its dimension by \Cref{prop:analytic-tropical-dimension}; apply
\Cref{thm:ch14-support-dimension}.
\end{proof}

\section{Pushforward and the projection formula}
\label{sec:ch14-pushforward}

Let $g:Y\to X$ be topologically proper.  For a current $T$ on $Y$, define
\[
                  \langle g_*T,\alpha\rangle
                    =\langle T,g^*\alpha\rangle.
\]
It suffices more generally that $\Supp(T)$ be proper over $X$.

\begin{proposition}
\label{prop:current-pushforward-formal}
Proper pushforward is weakly continuous, commutes with $J$, $\dprime$, and
$\ddprime$, and satisfies
\[
              g_*(T\wedge g^*\alpha)=g_*T\wedge\alpha
\]
for every smooth form $\alpha$.
\end{proposition}
\begin{proof}
Every formula follows by evaluation on a compactly supported test form and
functoriality of pullback.  Properness ensures that the pulled-back test form
meets $\Supp(T)$ in a compact set.

In detail.  \emph{$g_*T$ is a current.}  By
\Cref{prop:cld-form-pullback} the map $g^*$ is continuous
$A_X^{r,s}\to g_*A_Y^{r,s}$, and topological properness makes it carry
$A_{X,c}^{r,s}$ into $A_{Y,c}^{r,s}$ continuously: the preimage of a compact
set is compact, and a chart datum on $X$ pulls back to one on $Y$ with the
same seminorm estimates, because $g^*$ of a chart representative is the
representative for the composed moment map.  Hence the transpose is defined,
and it preserves relative order.  \emph{Weak continuity} is immediate, the
weak topology being that of pointwise convergence and $g_*$ being defined by
composition.

\emph{Commutation with $J$.}  $g^*J=Jg^*$ on forms
(\Cref{prop:cld-form-pullback}), so
$\langle g_*JT,\alpha\rangle=\langle JT,g^*\alpha\rangle
=\langle T,J^{-1}g^*\alpha\rangle=\langle T,g^*J^{-1}\alpha\rangle
=\langle Jg_*T,\alpha\rangle$.  In particular the pushforward of a symmetric
current is symmetric.

\emph{Commutation with $\dprime$.}  Let $T$ have bidegree $(u,v)$ on $Y$;
then $g_*T$ has the same bidegree if $\dim Y=\dim X$, and in general the
bidimension is preserved, so the sign $(-1)^{u+v+1}$ of
\eqref{eq:ch14-dprime-sign} is the same on both sides by the parity remark
after \Cref{lem:ch14-module-leibniz}.  Since $g^*\dprime=\dprime g^*$,
\[
 \langle g_*\dprime T,\alpha\rangle
 =\pm\langle T,g^*\dprime\alpha\rangle
 =\pm\langle T,\dprime g^*\alpha\rangle
 =\langle \dprime g_*T,\alpha\rangle
\]
with the same sign throughout.  The case of $\ddprime$ follows by conjugating
with $J$.

\emph{Projection formula.}  For a test form $\beta$,
\[
 \langle g_*(T\wedge g^*\alpha),\beta\rangle
 =\langle T\wedge g^*\alpha,g^*\beta\rangle
 =\langle T,g^*\alpha\wedge g^*\beta\rangle
 =\langle T,g^*(\alpha\wedge\beta)\rangle
 =\langle g_*T\wedge\alpha,\beta\rangle,
\]
using multiplicativity of $g^*$.  This is the current-level form of
\Cref{cor:ch11-projection-formula}.

\emph{Relaxing properness.}  If $\Supp T$ is $g$-compact, meaning that
$\Supp T\cap g^{-1}(C)$ is compact for every compact $C\subset X$, then for
$\alpha$ compactly supported the intersection
$\Supp T\cap\Supp g^*\alpha$ is compact, and
\Cref{lem:ch14-phi-support} gives a meaning to
$\langle T,g^*\alpha\rangle$.  All the formulas above persist verbatim, with
the same proofs, since each is an identity between pairings against forms
whose supports meet $\Supp T$ compactly.
\end{proof}

\begin{corollary}
\label{cor:ch14-pushforward-functorial}
If $h:Z\to Y$ and $g:Y\to X$ are topologically proper, then
$(g\circ h)_*=g_*h_*$ on currents.  Pushforward preserves each of the three
positivity types, and it preserves relative order.
\end{corollary}

\begin{proof}
Functoriality is the transpose of $(g\circ h)^*=h^*g^*$,
\Cref{prop:cld-form-pullback}.  Positivity is the third paragraph of the
proof of \Cref{prop:current-positivity-calculus}; order is the first
paragraph of the proof above.
\end{proof}

\subsection{Why there is no pullback of currents}
\label{sec:ch14-pullback-failure}

A general pullback of currents is not available: it would amount to
restricting a distribution to a possibly characteristic subspace.  We use
pullback only for smooth forms, continuous potentials before differentiating,
or in transverse/refined theories where it is separately constructed.

It is worth making the obstruction precise rather than leaving it as a
slogan, because the exact shape of the obstruction dictates the hypotheses of
\Cref{thm:cld-fibre-projection} and of the Bedford--Taylor products in
\Cref{chap:cld-bedford-taylor}.

\begin{proposition}
\label{prop:ch14-no-pullback}
Let $X$ be the tropical space $\mathbf R$ of level $1$, let $Y=\{0\}$ be the
tropical subspace consisting of the origin, and let $g:Y\hookrightarrow X$ be
the inclusion.  There is no map
\[
 g^*:D^{0,0}(X)\longrightarrow D^{0,0}(Y)
\]
that is weakly continuous and satisfies $g^*[u]=[u\circ g]$ for every smooth
function $u$ on $X$.
\end{proposition}

\begin{proof}
Note first that $D^{0,0}(Y)=\mathbf R$, since $A^{0,0}_{Y,c}=\mathbf R$;
under this identification $[c]=c$ for a constant $c$.

Fix $\chi\in C_c^\infty(\mathbf R)$ with $\chi(0)=1$ and
$\Supp\chi\subset[-1,1]$, and for $\varepsilon>0$ put
$u_\varepsilon(x)=\chi(x/\varepsilon)$, a smooth compactly supported function
on $X$.  For a test form $\alpha=a\,\dprime x\wedge\ddprime x$ with
$a\in C_c^\infty(\mathbf R)$ we have, by
\Cref{def:ch11-lattice-measure},
\[
 \bigl|\langle[u_\varepsilon],\alpha\rangle\bigr|
 =\Bigl|\int_{\mathbf R}\chi(x/\varepsilon)a(x)\,\rd x\Bigr|
 \le\|\chi\|_\infty\int_{-\varepsilon}^{\varepsilon}|a|
 \xrightarrow[\varepsilon\to0]{}0 .
\]
Hence $[u_\varepsilon]\to0$ weakly in $D^{0,0}(X)$.  On the other hand
\[
 g^*[u_\varepsilon]=[u_\varepsilon\circ g]=u_\varepsilon(0)=\chi(0)=1
\]
for every $\varepsilon$.  Weak continuity would give $0=1$.
\end{proof}

\begin{remark}
The mechanism is the one familiar from distribution theory: weak convergence
of currents controls integrals, not pointwise values, and a closed immersion
asks precisely for a pointwise value in the transverse directions.  Over
$\mathbf C$ the obstruction is quantified by wave front sets, and pullback is
available when the conormal of $Y$ avoids the wave front set of $T$.  No
such microlocal theory is available here, and none of the results of this
book requires one; the three substitutes of
\Cref{rem:ch14-pullback-survives} suffice.
\end{remark}

\begin{remark}[What survives]
\label{rem:ch14-pullback-survives}
Three substitutes are used in this book, and they should be kept apart.
\begin{enumerate}
\item Pullback along an \emph{affine isomorphism} of a real vector space,
  \Cref{def:ch14-affine-pullback}: no hypothesis needed, but no new
  information either.
\item Pullback of a \emph{potential}: if $u$ is a continuous function on $X$,
  then $g^*u$ is a continuous function on $Y$ and one \emph{defines}
  $g^*(\dprimeddprime[u]):=\dprimeddprime[g^*u]$.  This is not a pullback of
  currents; it depends on the potential, and two potentials with the same
  curvature current on $X$ may have different curvature currents on $Y$.  It
  is the mechanism of \Cref{thm:cld-fibre-projection}.
\item Transverse pullback: when $g$ is transverse to a cycle, the tropical
  intersection product of \Cref{thm:ch11-transverse-current-product} provides
  a pullback of the integration current with the correct multiplicity.  This
  is the route taken in \Cref{chap:delta-forms}.
\end{enumerate}
\end{remark}

\section{Integration currents and analytic cycles}
\label{sec:ch14-cycles}

Let $i:Y\hookrightarrow X$ be a closed analytic subspace, pure of dimension
$r$.  Define
\[
                   \langle\delta_Y,\alpha\rangle
                         =\int_Yi^*\alpha,
              \qquad \alpha\in A_{X,c}^{r,r}.
\]
This is a current of bidimension $(r,r)$.  If
$Z=\sum m_jY_j$ is a locally finite analytic $r$-cycle, put
$\delta_Z=\sum m_j\delta_{Y_j}$.

Three verifications are needed before this definition may be used, and each
is short with the tools already in place.

\begin{proposition}
\label{prop:ch14-delta-Y-current}
Let $Y\subset X$ be closed analytic, pure of dimension $r$, and let
$i:Y\to X$ be the inclusion.
\begin{enumerate}
\item The functional $\alpha\mapsto\int_Yi^*\alpha$ is a well-defined,
  strongly positive, symmetric current of bidimension $(r,r)$, of order $0$
  relative to every chart datum, and $\delta_Y$ is the pushforward along $i$
  of the integration current of the tropical space $Y$ itself.
\item For a locally finite $r$-cycle $Z=\sum_jm_jY_j$ the sum
  $\sum_jm_j\delta_{Y_j}$ is locally finite and defines a current; it is
  strongly positive when all $m_j\ge0$.
\item $\Supp\delta_Y=Y$ when $Y$ is reduced.
\end{enumerate}
\end{proposition}

\begin{proof}
(1) The integral $\int_Yi^*\alpha$ is defined by
\Cref{thm:cld-analytic-integration}, applied on the pure $r$-dimensional
space $Y$ with its canonical calibration, because $i^*\alpha$ is a smooth
$(r,r)$-form on $Y$ with compact support.  If $\alpha$ is weakly positive,
then $i^*\alpha$ is weakly positive by
\Cref{prop:ch09-positivity-inclusions}, and the integral of a weakly positive
top-degree form is nonnegative, weak positivity in top bidegree being
ordinary nonnegativity by \Cref{prop:ch09-positivity-inclusions}.  Symmetry
follows from $J$ acting trivially on top-degree forms.  Now
\Cref{prop:ch14-automatic-continuity}\,(2) applies and delivers continuity and
order $0$.  The identity $\delta_Y=i_*\delta_{Y}$ is
\Cref{sec:ch14-pushforward} applied to the proper map $i$, together with
$\langle\delta_Y,\beta\rangle=\int_Y\beta$ for $\beta$ on $Y$.

(2) Let $\alpha$ have compact support $K$.  By local finiteness only finitely
many $Y_j$ meet $K$, and for the others $i_j^*\alpha=0$; so the sum is
finite.  Continuity on a fixed chart datum follows because the same finite
subfamily works for every form supported in $K$.

(3) Let $U$ be open and meet $Y$.  Choose $y\in U\cap Y$ at which the local
tropical dimension of $Y$ equals $r$; such a point exists by
\Cref{prop:analytic-tropical-dimension}, since $Y$ is pure $r$-dimensional.
Choose a tropical chart $f$ of $Y$ at $y$ whose tropicalization has dimension
$r$ near $f_{\trop}(y)$, extend it to a chart of $X$, and set
$\alpha=\rho\,f_{\trop}^*\eta^{r}$ with $\eta=\dprimeddprime\|x\|^2$ and
$\rho\ge0$ a smooth cutoff supported in $U$ with $\rho(y)=1$.  Then
$i^*\alpha$ is a positive $(r,r)$-form on $Y$ with compact support, and it is
not identically zero, its coefficient at $y$ against the calibrated top
orientation being positive by \Cref{lem:ch14-eta-dominates}.  Hence
$\int_Yi^*\alpha>0$ by \cite[Corollary~7.3.7]{CLD25}, and $\delta_Y$ does not
vanish on $U$.
\end{proof}

\begin{example}[The current of integration over a skeleton]
\label{ex:ch14-skeleton-current}
Let $K$ be nontrivially valued and let $E=\mathbf G_{m,K}^{\An}/q^{\mathbf Z}$
be a Tate curve with $|q|<1$, as in \Cref{ex:trop-tate-curve}.  Its canonical
skeleton is the circle
\[
 \Sigma=\mathbf R/\ell\mathbf Z,
 \qquad \ell=-\log|q|=v_K(q)>0,
\]
with the retraction $\tau:E^{\An}\to\Sigma$ and the calibration given by the
constant lattice $\mathbf Z$.  The current of integration over $\Sigma$ is
\[
 \langle\delta_\Sigma,f\rangle=\int_\Sigma f|_\Sigma\,\rd x,
 \qquad f\in A_{E,c}^{0,0},
\]
where $\rd x$ is the lattice measure of \Cref{def:ch11-lattice-measure}.  It
is a positive current of bidimension $(1,1)$ on the $1$-dimensional space
$E^{\An}$, hence of bidegree $(0,0)$; its total mass is
\[
 \langle\delta_\Sigma,1\rangle=\ell=v_K(q).
\]
Moreover $\delta_\Sigma=\delta_{E}$, the integration current of the whole
curve.  Indeed both are strongly positive currents of bidimension $(1,1)$, and
by \Cref{thm:analytic-is-tropical-space} the canonical calibration of
$E^{\An}$ is carried by $\Sigma$ with multiplicity $1$: a test $(1,1)$-form
on $E^{\An}$ integrates to the integral of its restriction to $\Sigma$.
Equivalently, $\delta_\Sigma=\tau^*$ in the sense that
$\langle\delta_\Sigma,f\rangle=\langle\delta_{\Sigma},f\circ\tau\rangle$
whenever $f$ is $\tau$-invariant, which is the content of
\Cref{cor:trop-skeleton-intrinsic}.

The example is the smallest one in which the calibration is a genuine
invariant: the mass $v_K(q)$ is the valuation of the Tate parameter, and it
changes if $q$ is replaced by $q^2$, even though the underlying topological
space does not change.
\end{example}

\begin{example}[Cycle currents on the Berkovich line]
\label{ex:ch14-line-cycles}
Return to $X=\mathbf P^{1,\An}_k$ with $k$ trivially valued and
algebraically closed, as in \Cref{sec:ch14-line-examples}.  The analytic
cycles of dimension $0$ on $X$ are the locally finite formal sums
$Z=\sum_bm_b\,b$ over rigid points, and
\[
 \delta_Z=\sum_bm_b\,\delta_{x_b}
\]
is the corresponding measure; the sum is finite against any test function by
\Cref{lem:ch14-star-smooth}.  The only cycle of dimension $1$ is
$m\cdot X$, with $\delta_{mX}=m\,\delta_X$; by
\Cref{cor:ch14-delta-x-closed} it is closed, and its ``mass'' against the
constant test form is the total length of the skeleton of $X$, which is zero
because the skeleton of $\mathbf P^{1,\An}_k$ over a trivially valued field is
a point.  This is the degenerate case in which the calibration carries no
metric information; contrast \Cref{ex:ch14-skeleton-current}.

Finally, \Cref{ex:ch14-line-model-function} exhibits a difference of two cycle
currents as a $\dprimeddprime$: with $r=1$ and $c_1=1$, the identity there
reads
\[
 \dprimeddprime[u]=\delta_{p_1}-\delta_\eta,
\]
with $u=\log\max\{|T-b_1|,\rme^{-1}\}$ on the branch $b_1$ and $u=0$
elsewhere.  This is the truncated local model of the Poincar\'e--Lelong
formula \eqref{eq:appJ-cld-pl}; the untruncated statement, and the passage
from a truncation to a divisor, are the subject of
\Cref{chap:cld-poincare-lelong}.
\end{example}

\begin{theorem}
\label{thm:current-fundamental-cycle}
For every pure-dimensional closed analytic subspace $Y\subset X$,
\[
                         \delta_Y=\delta_{[Y]}.
\]
Thus nilpotent and embedded structure enters precisely through the generic
multiplicities of the fundamental cycle.
\end{theorem}
\begin{proof}
Write $Y_1,\ldots,Y_s$ for the irreducible components of
$Y_{\mathrm{red}}$ and let
$m_j=\operatorname{length}_{\mathcal O_{Y,\eta_j}}
\mathcal O_{Y,\eta_j}$ be the generic multiplicity along $Y_j$.  Remove the
union $B$ of the proper intersections $Y_i\cap Y_j$ and the loci where an
irreducible component is not generically reduced.  This is a closed analytic
set of dimension less than $r$.  A test form of bidegree $(r,r)$ restricts
to zero on $B$ by dimension, so deleting $B$ does not change either
integration current.

On $Y\setminus B$ the components are disjoint.  Choose a tropical chart at a
point of $Y_j\setminus B$.  The canonical calibration of the nonreduced
analytic space is defined by the initial cycle of its affinoid algebra.
Localization at the generic prime of $Y_j$ shows that this initial cycle is
$m_j$ times the initial cycle of $(Y_j)_{\mathrm{red}}$: multiplicity is the
length of the localized algebra, and length is unchanged by the graded
reduction used to define the calibration.  Hence chartwise integration over
$Y$ equals $m_j$ times integration over $Y_j$.

Summing over the disjoint components gives
\[
          \delta_Y=\sum_jm_j\delta_{Y_j}=\delta_{[Y]}.
\]
The removal and graded-length arguments are
\cite[Proposition~16.6.3]{CLD25}.  They also show why replacing $Y$ by its
reduction without the coefficients $m_j$ would be incorrect.
\end{proof}

\begin{proposition}
\label{prop:current-cycle-closed}
If $X$ is boundaryless, the integration current $\delta_X$ is
$\dprime$- and $\ddprime$-closed.  More generally, integration currents of
analytic cycles are closed.
\end{proposition}
\begin{proof}
For $\delta_X$ this is Stokes.  For a cycle, apply Stokes on each pure
component and sum locally finitely.  Equivalently use proper pushforward from
each component and \Cref{prop:current-pushforward-formal}.
\end{proof}

\section{Fibre integration}
\label{sec:ch14-fibre-integration}

Fibre integration is subtler than formal pushforward because it asserts that
the pushforward of a form is represented by a function.  Let
$p:Y\to X$ be a morphism of pure-dimensional topologically separated spaces,
$\dim X=m$, $\dim Y=n$, and $d=n-m$.  Let $\omega$ be a $(d,d)$-form on $Y$
whose support is proper over $X$.  On a fibre of pure dimension $d$, set
\[
                         u_\omega(x)=\int_{Y_x}\omega;
\]
set it equal to zero when the fibre does not have dimension $d$.

\begin{theorem}
\label{thm:cld-fibre-integration}
Assume that the flat locus of $p$ is dense in $Y$.  If $\omega$ is a
$G$-form of bidegree $(d,d)$ with locally bounded tropically Borel
coefficients and $p$-compact support, then $u_\omega$ is locally integrable
and
\[
              \int_Y\omega\wedge p^*\alpha
                =\int_Xu_\omega\,\alpha
\]
for every compactly supported smooth $(m,m)$-form $\alpha$.  Equivalently,
\[
                             p_*[\omega]=[u_\omega].
\]
\end{theorem}
\begin{proof}
Work first over a relatively compact tropical chart on $X$.  Adjoin to the
pullbacks of its moment functions finitely many invertible analytic functions
on $Y$ that represent $\omega$.  This gives a product moment map on $Y$ and
an integer-affine projection $F$ from its tropicalization to the base
tropicalization.  Subdivide both complexes so that $F$ is affine of constant
rank on every cell and every coefficient of $\omega$ is tropically Borel on
the cells.

On the dense flat locus, a top cell of $Y$ mapping to a top cell of $X$ has
relative dimension $d$.  In lattice coordinates adapted to $F$, ordinary
Fubini's theorem gives
\[
 \int_\sigma\omega\wedge F^*\alpha
   =\int_{F(\sigma)}
       \left(\int_{\sigma_x}\omega\right)\alpha(x).
\]
If the adapted coordinates are not unimodular, change of variables
introduces the index
$[N_{F(\sigma)}:F(N_\sigma/N_{\sigma_x})]$.  The calibration pushforward
formula \eqref{eq:analytic-tropical-pushforward} says that the same index,
together with the residue-field degree, is exactly the ratio between the
source and target analytic multiplicities.  Thus the displayed equality is
valid with the canonical calibrations, not only with Euclidean volume.

Cells on which the rank drops have image of lower dimension and contribute
zero to a top-degree test.  The complement of the flat locus has no
top-dimensional component by density; its calibrated contribution is
therefore also zero.  Summing the cellwise identity proves the formula for a
form supported in the chosen charts.  A partition of unity and common
product refinements prove it globally.  Local boundedness of the tropical
coefficients bounds the fibre integrals by finitely many polyhedral volumes
on compact subsets, so $u_\omega$ is locally integrable.  Approximation by
simple Borel coefficients and dominated convergence remove the temporary
cellwise-constant assumption.

This is the calibrated Fubini argument of
\cite[Theorem~16.5.4 and Corollary~16.5.5]{CLD25}.  Dense flatness is used in
the dimension and degree step; omitting it would lose the multiplicities of
nonreduced fibres.
\end{proof}

\begin{theorem}
\label{thm:cld-fibre-projection}
Under the hypotheses of \Cref{thm:cld-fibre-integration}, let
$u_1,\ldots,u_r$ be tropically differences of convex functions on $X$, let
$\omega$ be smooth of bidegree $(d,d)$ with $p$-compact support, and let
$\alpha$ be a compactly supported $(m-r,m-r)$-form.  Then
\[
\begin{split}
 &\left\langle
   \bigwedge_{j=1}^r\dprimeddprime(p^*u_j),
   \omega\wedge p^*\alpha\right\rangle\\
 &\hspace{25mm}=
 \left\langle\bigwedge_{j=1}^r\dprimeddprime u_j,
                     u_\omega\alpha\right\rangle .
\end{split}
\]
\end{theorem}
\begin{proof}
Assume first that the $u_j$ are smooth.  Pullback commutes with both
differentials, hence
\[
 \dprimeddprime(p^*u_j)=p^*(\dprimeddprime u_j).
\]
Apply \Cref{thm:cld-fibre-integration} with
$\bigwedge_j\dprimeddprime u_j\wedge\alpha$ in place of the base test form.
The resulting identity is exactly the formula in the statement.

On a tropical chart, a continuous convex function is uniformly approximated
on relatively compact subsets by smooth convex functions.  A tropically
difference-of-convex function is the difference of two such approximations.
Choose approximations $u_{j,k}$ simultaneously on the finitely many charts
meeting the supports.  The Bedford--Taylor products
$\bigwedge_j\dprimeddprime u_{j,k}$ converge weakly to the corresponding
products for $u_j$, and the same is true after pullback to $Y$.  The fibre
function $u_\omega$ is locally integrable by
\Cref{thm:cld-fibre-integration}; the uniform local bounds in the
Bedford--Taylor construction allow it to be used as the coefficient of the
test form.  Passing to the limit in the smooth identity gives the displayed
formula.  Independence of the decompositions into differences follows from
the same continuity.  These approximation and convergence steps are
\cite[Theorem~16.5.7]{CLD25}; the local products themselves are constructed
in \Cref{chap:cld-bedford-taylor}.
\end{proof}

\subsection{Two fibre integrations carried to numbers}
\label{sec:ch14-fibre-numbers}

The hypotheses of \Cref{thm:cld-fibre-integration} are opaque until one has
run the computation once with a degree in it.  We do so twice: first with a
unimodular projection, where nothing happens, and then with a doubling map,
where the calibrated degree $2$ of the fibre appears on the nose.  In both
examples $K$ is nontrivially valued, and we use
\Cref{prop:trop-torus-skeleton} to identify the tropicalization of an
analytic torus with $\mathbf R^{\dim}$, calibrated with multiplicity $1$, and
\Cref{thm:cld-analytic-integration} to compute integrals there.

Throughout, $\rho\in C_c^\infty(\mathbf R)$ is a fixed nonnegative bump with
\[
 \int_{\mathbf R}\rho(t)\,\rd t=1,
 \qquad \Supp\rho\subset[-1,1].
\]

\begin{example}[A unimodular projection]
\label{ex:ch14-fibre-projection-numbers}
Let $Y=(\mathbf G_{m,K}^2)^{\An}$ with tropical coordinates $(t_1,t_2)$, let
$X=(\mathbf G_{m,K})^{\An}$ with tropical coordinate $s$, and let
$p:Y\to X$ be given by $p^\#S=T_1$; tropically $F(t_1,t_2)=t_1$.  Here
$n=2$, $m=1$, $d=1$, and $p$ is flat everywhere, so the density hypothesis of
\Cref{thm:cld-fibre-integration} holds.  Take
\[
 \omega=\trop^*\bigl(\rho(t_1)\rho(t_2)\,\dprime t_2\wedge\ddprime t_2\bigr),
\]
a smooth $(1,1)$-form on $Y$ with compact support, hence $p$-compact support.

The fibre $Y_x$ over a point $x$ with $s(x)=\sigma$ is a torus of dimension
$1$ over $\mathcal H(x)$ with tropical coordinate $t_2$ and multiplicity $1$,
so
\[
 u_\omega(x)=\int_{Y_x}\omega
 =\rho(\sigma)\int_{\mathbf R}\rho(t_2)\,\rd t_2
 =\rho(\sigma).
\]
Now let $\alpha=\trop^*\bigl(a(s)\,\dprime s\wedge\ddprime s\bigr)$ with
$a\in C_c^\infty(\mathbf R)$.  Since $p^\#S=T_1$ we have
$p^*\alpha=\trop^*(a(t_1)\dprime t_1\wedge\ddprime t_1)$, and since
$\dprime t_2\wedge\ddprime t_2$ has even total degree it commutes past
$\dprime t_1\wedge\ddprime t_1$, so
\[
 \omega\wedge p^*\alpha
 =\rho(t_1)\rho(t_2)a(t_1)\,\tau_2 .
\]
Therefore
\[
 \int_Y\omega\wedge p^*\alpha
 =\int_{\mathbf R^2}\rho(t_1)\rho(t_2)a(t_1)\,\rd t_1\rd t_2
 =\int_{\mathbf R}\rho(t_1)a(t_1)\,\rd t_1
 =\int_Xu_\omega\,\alpha .
\]
With $a\equiv1$ near $[-1,1]$ both sides equal $1$.  No index appears, because
$F$ is the projection $\mathbf Z^2\to\mathbf Z$ onto a direct summand.
\end{example}

\begin{example}[A doubling map: the degree appears]
\label{ex:ch14-fibre-degree-numbers}
Keep $Y$ and $X$ as above but let $q:Y\to X$ be given by
$q^\#S=T_1^2$; tropically $G(t_1,t_2)=2t_1$.  Now
$G(\mathbf Z^2)=2\mathbf Z\subset\mathbf Z$, of index $2$, and analytically
the fibre
\[
 Y_x=\{T_1^2=S(x)\}\times\mathbf G_m
\]
is finite of degree $2$ over a $1$-dimensional torus: the algebra
$\mathcal H(x)[T_1]/(T_1^2-S(x))$ is a degree-$2$ \'etale
$\mathcal H(x)$-algebra as soon as $\operatorname{char}\mathcal H(x)\neq2$,
and in any case its length is $2$.  Hence the calibrated degree of $Y_x$ is
$2$, by \Cref{thm:trop-degree-graded} and \Cref{def:trop-calibrated-degree}.

Take the same $\omega$ as before.  On the fibre over $s(x)=\sigma$ we have
$t_1=\sigma/2$ on each of the two sheets, so
\[
 u_\omega(x)=\int_{Y_x}\omega
 =2\,\rho(\sigma/2)\int_{\mathbf R}\rho(t_2)\,\rd t_2
 =2\,\rho(\sigma/2),
\]
the factor $2$ being the calibrated degree.  On the other side,
$q^\#S=T_1^2$ gives $\dprime s=2\dprime t_1$ and $\ddprime s=2\ddprime t_1$,
hence
\[
 q^*\alpha=4\,a(2t_1)\,\trop^*\bigl(\dprime t_1\wedge\ddprime t_1\bigr),
\]
and
\[
 \int_Y\omega\wedge q^*\alpha
 =4\int_{\mathbf R^2}\rho(t_1)\rho(t_2)a(2t_1)\,\rd t_1\rd t_2
 =4\int_{\mathbf R}\rho(t_1)a(2t_1)\,\rd t_1 .
\]
Substituting $\sigma=2t_1$, so $\rd t_1=\rd\sigma/2$,
\[
 \int_Y\omega\wedge q^*\alpha
 =2\int_{\mathbf R}\rho(\sigma/2)a(\sigma)\,\rd\sigma
 =\int_{\mathbf R}u_\omega(\sigma)a(\sigma)\,\rd\sigma
 =\int_Xu_\omega\,\alpha,
\]
as \Cref{thm:cld-fibre-integration} predicts.  With $a\equiv1$ on
$[-2,2]$ both sides equal
\[
 2\int_{\mathbf R}\rho(\sigma/2)\,\rd\sigma=4 .
\]
Thus the total mass is $4$, not $1$: a factor $2$ from the calibrated degree
of the fibre and a factor $2$ from the Jacobian of $\sigma\mapsto\sigma/2$ on
the base.  Had one omitted the calibrated degree and simply written
$u_\omega=\rho(\sigma/2)$, the two sides would have differed by a factor
$2$; this is the numerical form of the multiplicity assertion inside the
proof of \Cref{thm:cld-fibre-integration}, and it is the check that detects
the commonest error.
\end{example}

\begin{remark}[Where each hypothesis is used]
\label{rem:ch14-fibre-hypotheses}
\Cref{thm:cld-fibre-integration} has three hypotheses, and each fails
detectably.
\begin{enumerate}
\item \emph{$p$-compact support.}  Drop it and take
  $\omega=\trop^*(\rho(t_1)\dprime t_2\wedge\ddprime t_2)$ in
  \Cref{ex:ch14-fibre-projection-numbers}: the fibre integral is
  $+\infty$ and $u_\omega$ is not defined.  This is the fibre-integration
  analogue of \Cref{ex:ch11-pushforward-nonproper}.
\item \emph{Dense flat locus.}  Let $X$ be nonreduced and $p$ the immersion
  $X_{\mathrm{red}}\hookrightarrow X$, which is not flat.  Then $d=0$ and
  $u_\omega=1$ on $X_{\mathrm{red}}$, while $p_*[\omega]$ carries the generic
  multiplicity of $X$; the identity fails by exactly that multiplicity, as
  \Cref{thm:current-fundamental-cycle} makes precise.
\item \emph{Locally bounded tropically Borel coefficients.}  These are what
  make the cellwise Fubini step legitimate; without a local bound the fibre
  integral need not be finite even when the support is $p$-compact.
\end{enumerate}
\end{remark}

\section{Support and compactness principles}
\label{sec:ch14-compactness}

The full support theory of positive closed currents uses scalar extension and
skeletal descriptions.  We record the conclusions needed later and diagnose
the Hausdorff index in the scalar-extension case.  The archival comments in
the next two inputs are retained solely so that the bytes used by the book
remain identical to the production-verified packet.  For the same reason,
the theorem, index remark, and diagnostic proposition are grouped before
their two proofs; each proof is explicitly headed by the result it proves.

% Quarantine-only statement.
% Date: 29 August 2026.
% This file is not a live-book input and makes no novelty or priority claim.

\begin{theorem}
\label{thm:ch14-cld-support-corrected-packet}
Let $K$ be a complete non-Archimedean field, and let $X$ be a
topologically separated, boundaryless, pure $n$-dimensional $K$-analytic
space.
\begin{enumerate}
\item\label{itm:ch14-cld-support-corrected-nontrivial}
Assume that $|K^\times|$ is nontrivial.  Let $0\leq p\leq n$, let $T$ be a
positive closed current of bidimension $(p,p)$ on $X$, and let $A\subset X$
be closed and tropically $H^p$-negligible.  Thus, for every compact analytic
domain $V$ of $X$ and every tropical chart
$f_{\trop}:V\to\mathbf R^N$, the set $f_{\trop}(A\cap V)$ has zero
$p$-dimensional Hausdorff measure.  Then, for every smooth
$\alpha\in A_X^{p,p}(X)$, the signed Radon measure $T\wedge\alpha$ satisfies
\begin{equation}
\label{eq:ch14-cld-support-corrected-nontrivial}
                       |T\wedge\alpha|(A)=0.
\end{equation}

\item\label{itm:ch14-cld-support-corrected-trivial}
Assume that $K$ is trivially valued.  Let $0\leq q\leq n$, let
$u_1,\ldots,u_q$ be continuous dpsh functions on $X$, and put
\begin{equation}
\label{eq:ch14-cld-support-corrected-bt-current}
                    T=\bigwedge_{j=1}^{q}\dprimeddprime u_j.
\end{equation}
If $A\subset X$ is closed and tropically $H^{n-q}$-negligible, then, for
every smooth form
\[
                         \alpha\in A_X^{n-q,n-q}(X),
\]
\begin{equation}
\label{eq:ch14-cld-support-corrected-trivial}
                         |T\wedge\alpha|(A)=0.
\end{equation}
Here the empty product for $q=0$ is the integration current of $X$.
\end{enumerate}
\end{theorem}

\begin{remark}
\label{rem:ch14-cld-support-corrected-index}
In \cref{itm:ch14-cld-support-corrected-nontrivial}, the integer $p$ is a
\emph{bidimension}.  In \cref{itm:ch14-cld-support-corrected-trivial}, the
integer $q$ counts curvature factors.  The current in
\eqref{eq:ch14-cld-support-corrected-bt-current} has bidegree $(q,q)$ and
bidimension $(n-q,n-q)$.  Hence the exponent forced by
\cref{itm:ch14-cld-support-corrected-nontrivial} is $n-q$.  The exponent
$H^q$ printed in \cite[Corollary~16.8.6]{CLD25} is therefore an index error.
\end{remark}

\begin{proposition}
\label{prop:ch14-cld-support-printed-exponent-counterexample}
Let $K$ be an algebraically closed complete trivially valued field, let
$X=(\mathbf G_{m,K})^{\mathrm{an}}$, let $z$ be its coordinate, and put
\[
               \tau=\log|z|,
               \qquad
               u=\max\{0,\tau\}.
\]
If $\eta_1$ denotes the unit Gauss point, then $u$ is a continuous CLD-psh
function and
\begin{equation}
\label{eq:ch14-cld-support-kink-dirac}
                         \dprimeddprime[u]=\delta_{\eta_1}.
\end{equation}
The singleton $A=\{\eta_1\}$ is tropically $H^1$-negligible.  Consequently,
when $n=q=1$, the assertion obtained from
\cref{itm:ch14-cld-support-corrected-trivial} by replacing $H^{n-q}$ with
$H^q$ is false.
\end{proposition}

% Quarantine-only proof.
% Date: 29 August 2026.
% This file is not a live-book input and makes no novelty or priority claim.

\begin{proof}[Proof of
\cref{thm:ch14-cld-support-corrected-packet}]
We first prove
\cref{itm:ch14-cld-support-corrected-nontrivial}.  This is
\cite[Theorem~16.8.5]{CLD25}; we give the argument because the proper
tropical pushforward is essential and because it fixes the exponent without
any appeal to a ball-growth estimate on $X$.

The assertion is local.  Fix $x\in X$.  After shrinking, we may suppose
that $X$ is paracompact and that $\alpha$ is tropical.  We may also reduce to
the case in which $\alpha$ is positive.  Indeed, $T$ is symmetric, so its
wedge with the antisymmetric part of $\alpha$ is zero.  On a relatively
compact tropical chart, the symmetric form $\alpha$ is a difference of two
positive forms: if $\Theta$ is a smooth strictly positive $(p,p)$-form on a
neighborhood of the relevant compact set, then
$C\Theta+\alpha$ and $C\Theta$ are positive for $C\gg1$.  A finite smooth
partition of unity gives the same decomposition on an arbitrary compact
subset.  It is therefore enough to prove the assertion for each of the two
positive measures separately.

Choose a compact analytic neighborhood $V$ of $x$ and a moment map
\[
                         f:V\longrightarrow\mathbf G_m^N
\]
such that
\begin{equation}
\label{eq:ch14-cld-support-boundary-separation}
                 f_{\trop}(x)\notin f_{\trop}(\partial V)
\end{equation}
and such that $f$ tropicalizes $\alpha|_V$.  The separating moment exists
by \cite[Lemma~12.3.6]{CLD25}; concatenating it with a moment tropicalizing
$\alpha$ preserves \eqref{eq:ch14-cld-support-boundary-separation}.  Thus
there is a smooth $(p,p)$-superform $\omega$, defined near
$f_{\trop}(V)$, for which
\begin{equation}
\label{eq:ch14-cld-support-alpha-tropical}
                         \alpha|_V=f_{\trop}^*\omega.
\end{equation}
Multiplication by a plateau function permits us to regard $\omega$ as
compactly supported on $\mathbf R^N$ without changing it near the compact
set used below.

Set
\begin{equation}
\label{eq:ch14-cld-support-omega-U}
 \Omega=\mathbf R^N\setminus f_{\trop}(\partial V),
 \qquad
 U=f_{\trop}^{-1}(\Omega)\subset V.
\end{equation}
The compactness of $\partial V$ makes $\Omega$ open.  The set $U$ is
disjoint from $\partial V$, hence lies in $\operatorname{Int}(V/X)$ and is
open in $X$; by \eqref{eq:ch14-cld-support-boundary-separation}, it contains
$x$.  Moreover,
\begin{equation}
\label{eq:ch14-cld-support-proper-tropicalization}
                     f_{\trop}|_U:U\longrightarrow\Omega
\end{equation}
is proper: it is obtained from the proper map from the compact space $V$ to
$\mathbf R^N$ by restriction over $\Omega$.

Properness makes the following tropical pushforward well-defined:
\begin{equation}
\label{eq:ch14-cld-support-pushed-current}
 \langle T_\Omega,\beta\rangle
   :=\langle T|_U,f_{\trop}^*\beta\rangle,
 \qquad \beta\in A_c^{p,p}(\Omega).
\end{equation}
Pullback preserves the relevant positivity cone and commutes with
$\dprime$ and $\ddprime$.  Hence $T_\Omega$ is a positive closed
supercurrent of bidimension $(p,p)$ on $\Omega$.  In Lagerberg's bidegree
notation it has bidegree $(N-p,N-p)$.  Lagerberg's Euclidean support
theorem, \cite[Corollary~4.1]{Lag12Super}, says that the coefficient measures
of such a current vanish on every set of zero $p$-dimensional Hausdorff
measure.  (The same result is Corollary~3.5 in the arXiv version.)

The set
\[
                       E=f_{\trop}(A\cap U)
\]
has zero $p$-dimensional Hausdorff measure: it is contained in
$f_{\trop}(A\cap V)$, which is null by hypothesis.  It follows that all
coefficient measures of $T_\Omega$ vanish on $E$, and therefore the Radon
measure $T_\Omega\wedge\omega$ does not charge $E$.

For completeness, we now transfer this coefficientwise statement back to
$U$ without a cancellation argument.  We are treating the case in which
$\alpha$ is positive, so
\[
                           \mu=T\wedge\alpha
\]
is a positive Radon measure.  By
\eqref{eq:ch14-cld-support-alpha-tropical} and the definition
\eqref{eq:ch14-cld-support-pushed-current}, its ordinary measure-theoretic
pushforward is
\begin{equation}
\label{eq:ch14-cld-support-measure-pushforward}
                  (f_{\trop}|_U)_*(\mu|_U)
                         =T_\Omega\wedge\omega.
\end{equation}
Indeed, both sides take a compactly supported smooth function $h$ on
$\Omega$ to
$\langle T|_U,f_{\trop}^*(h\omega)\rangle$; equality for Radon measures
then follows from density of smooth compactly supported functions in
continuous compactly supported functions.

Let $B\subset A\cap U$ be compact.  Its image $C=f_{\trop}(B)$ is compact
and contained in $E$.  From
\eqref{eq:ch14-cld-support-measure-pushforward} we obtain
\[
 0=(T_\Omega\wedge\omega)(C)
   =\mu\bigl(f_{\trop}^{-1}(C)\cap U\bigr)
   \geq\mu(B).
\]
Thus $\mu(B)=0$.  Inner regularity gives $\mu(A\cap U)=0$.  Repeating this
argument for the two positive terms in the decomposition of $\alpha$ proves
\eqref{eq:ch14-cld-support-corrected-nontrivial} near $x$, and locality
proves it on $X$.

We next prove
\cref{itm:ch14-cld-support-corrected-trivial}.  Put
\begin{equation}
\label{eq:ch14-cld-support-p-equals-n-minus-q}
                              p=n-q.
\end{equation}
By multilinearity of the Bedford--Taylor product, it is enough to treat the
case in which every $u_j$ is psh.  The current $T$ is then positive and
closed.  The dpsh case follows by writing each $u_j$ locally as a difference
of two continuous psh functions and expanding; all support, pushforward,
and null-set assertions below are stable under finite linear combinations.

Fix $0<r<1$ and form the one-radius Gauss extension
\[
              K_r=K((\varpi)),\qquad |\varpi|=r.
\]
Equivalently, $K_r=\mathscr H(\eta_r)$ for the radius-$r$ Gauss point.  The
radius is $K$-free because $K$ is trivially valued and $r\ne1$.  In
particular, $K_r$ is nontrivially valued and universally multiplicative;
see \cite[Example~11.1.3 and \S12.7.4]{CLD25}.  Let
\[
 \pi:X_r=X\widehat\otimes_KK_r\longrightarrow X,
 \qquad
 \sigma:X\longrightarrow X_r
\]
be the projection and its continuous Shilov section.  Pure dimension and
boundarylessness are preserved by this scalar extension
\cite[\S\S11.2.5 and 11.4.1]{CLD25}, so the nontrivially valued part already
proved applies to $X_r$.

We need the following exact consequence of the Gauss construction:
\begin{equation}
\label{eq:ch14-cld-support-shilov-null}
        A\text{ tropically }H^p\text{-negligible}
        \quad\Longrightarrow\quad
        \sigma(A)\text{ tropically }H^p\text{-negligible in }X_r.
\end{equation}
Here is the verification.  The issue is local, so we may replace $X$ by an
affinoid domain.  Then $A$ is compact and $\sigma(A)$ is compact, hence
closed.  Let $W$ be a compact analytic domain of $X_r$ and
$f=(f_1,\ldots,f_m)$ a moment on $W$.  Applied to each $f_i$, and then
refined simultaneously, \cite[Lemma~16.8.7]{CLD25} gives a finite compact
cover $(C_\lambda)$ of $\sigma^{-1}(W)$, base moments $g_\lambda$ defined on
neighborhoods of $C_\lambda$, and vectors $e_\lambda\in\mathbf Z^m$ such
that
\begin{equation}
\label{eq:ch14-cld-support-gauss-moment-comparison}
 f_{\trop}\circ\sigma
       =g_{\lambda,\trop}+e_\lambda\log r
                 \quad\text{on }C_\lambda.
\end{equation}
After covering each $C_\lambda$ by finitely many compact analytic domains
on which $g_\lambda$ is a tropical chart, the hypothesis on $A$ shows that
$g_{\lambda,\trop}(A\cap C_\lambda)$ is $H^p$-null.  Translation by the
constant vector $e_\lambda\log r$ preserves Hausdorff measure.  The finite
union furnished by
\eqref{eq:ch14-cld-support-gauss-moment-comparison} proves
\eqref{eq:ch14-cld-support-shilov-null}.

Write
\[
 u_{j,r}=\pi^*u_j,
 \qquad
 T_r=\bigwedge_{j=1}^q\dprimeddprime u_{j,r},
 \qquad
 \alpha_r=\pi^*\alpha.
\]
The current $T_r$ has bidegree $(q,q)$ and therefore bidimension
$(p,p)$ by \eqref{eq:ch14-cld-support-p-equals-n-minus-q}.  Applying
\cref{itm:ch14-cld-support-corrected-nontrivial} on $X_r$ to the
$H^p$-null set $\sigma(A)$ gives
\begin{equation}
\label{eq:ch14-cld-support-upstairs-null}
                         |T_r\wedge\alpha_r|(\sigma(A))=0.
\end{equation}

Two exact scalar-extension identities now finish the proof.  Since
$\alpha_r$ has bidegree $(p,p)=(n-q,n-q)$,
\cite[Proposition~16.8.2]{CLD25} gives
\begin{equation}
\label{eq:ch14-cld-support-shilov-support}
                \operatorname{supp}(T_r\wedge\alpha_r)
                                  \subseteq\sigma(X).
\end{equation}
The pushforward formula in \cite[Proposition~8.7.5]{CLD25}, in the scalar
extension form recalled immediately after
\cite[Proposition~16.8.2]{CLD25}, gives
\begin{equation}
\label{eq:ch14-cld-support-gauss-pushforward}
                  \pi_*(T_r\wedge\alpha_r)=T\wedge\alpha.
\end{equation}
Because $\pi\circ\sigma=\operatorname{id}_X$, one has
\[
        \pi^{-1}(A)\cap\sigma(X)=\sigma(A).
\]
Together with
\eqref{eq:ch14-cld-support-upstairs-null} and
\eqref{eq:ch14-cld-support-shilov-support}, this says that
$T_r\wedge\alpha_r$ does not charge $\pi^{-1}(A)$.  Pushing forward through
\eqref{eq:ch14-cld-support-gauss-pushforward} proves
\eqref{eq:ch14-cld-support-corrected-trivial}.  Notice that the Hausdorff
exponent used upstairs is
$p=n-q$, the bidimension of $T_r$, and not the number $q$ of curvature
factors.
\end{proof}

\begin{proof}[Proof of
\cref{prop:ch14-cld-support-printed-exponent-counterexample}]
For $t\in\mathbf R$, let $s(t)$ be the Gauss seminorm
\[
 \left|\sum_{j\in\mathbf Z}a_jz^j\right|_{s(t)}
       =\max_{a_j\ne0}\mathrm e^{jt}.
\]
Its image is the canonical skeleton of $X$, and
$\tau(s(t))=t$; in particular $s(0)=\eta_1$.  The function
$h(t)=\max\{0,t\}$ is continuous and convex.  Hence
$u=h\circ\tau$ is CLD-psh by the tropical convexity criterion
\cite[Proposition~8.2.5]{CLD25}.

We compute its curvature with the lattice coordinate $t$.  The
one-dimensional tropical current formula identifies the curvature of this
tropical potential with the distributional Hessian of $h$ on the canonical
skeleton.  Explicitly, if $\varphi$ is a compactly supported smooth test
function, the calibrated integration formula and two integrations by parts
give
\begin{align}
 \left\langle\dprimeddprime[u],\varphi\right\rangle
   &=\int_{\mathbf R}\max\{0,t\}
           \frac{\mathrm d^2}{\mathrm dt^2}
             \bigl(\varphi\circ s\bigr)(t)\,\mathrm dt \notag\\
   &=-\int_0^\infty
           \frac{\mathrm d}{\mathrm dt}
             \bigl(\varphi\circ s\bigr)(t)\,\mathrm dt
     =\varphi(s(0)).
\label{eq:ch14-cld-support-kink-test}
\end{align}
There is no boundary term because $\varphi$ has compact support.  Equivalently,
this is the slope jump $1-0=1$ at the only corner of $h$.  The tropical
Bedford--Taylor construction, or directly
\cite[Theorem~8.10.3]{CLD25}, shows that no curvature is present off this
zero-dimensional skeletal corner.  Thus
\eqref{eq:ch14-cld-support-kink-test} proves
\eqref{eq:ch14-cld-support-kink-dirac}.

Every tropical chart sends the singleton $A=\{\eta_1\}$ either to the empty
set or to a singleton.  Its one-dimensional Hausdorff measure is therefore
zero, so $A$ is tropically $H^1$-negligible.  On the other hand,
\[
                    \dprimeddprime[u](A)=1.
\]
For $n=q=1$ and $\alpha=1$, the exponent printed in
\cite[Corollary~16.8.6]{CLD25} is $H^q=H^1$ and would force this mass to be
zero.  This contradiction proves that the printed exponent cannot be
correct.  With the corrected exponent $H^{n-q}=H^0$, there is no
contradiction: a nonempty singleton has positive $H^0$-measure.
\end{proof}


\begin{remark}
\label{rem:ch14-support-route}
The first proof above is the proper-tropical-pushforward route of
\cite[Theorem~16.8.5]{CLD25}; its only Euclidean support input is
\cite[Corollary~4.1]{Lag12Super}.  It asserts no local Berkovich ball-growth
estimate.  The second proof uses the Gauss--Shilov mechanism of
\cite[Proposition~16.8.2 and Proposition~8.7.5]{CLD25} and
\cite[Lemma~16.8.7 and the proof of Corollary~16.8.6]{CLD25}, with the
Hausdorff exponent corrected as in
\cref{rem:ch14-cld-support-corrected-index}.
\end{remark}

\begin{proposition}
\label{prop:positive-current-compactness}
Let $(T_i)$ be positive currents of fixed bidegree.  Suppose that on every
relatively compact chart $U$ there is a smooth strictly positive test form
$\Omega_U$ of complementary bidegree and a constant $C_U$ such that
\[
                 \langle T_i,\chi\Omega_U\rangle\leq C_U
\]
for one cutoff $\chi$ equal to one on a prescribed compact subset of $U$.
Then the family is pointwise bounded on every fixed compact test-form space.
Consequently every weak cluster point is positive.
\end{proposition}
\begin{proof}
On fixed compact chart data, every symmetric test form $\alpha$ satisfies
\[
 -C\|\alpha\|_{C^0}\Omega_U\leq\alpha\leq
 C\|\alpha\|_{C^0}\Omega_U.
\]
Positivity gives a uniform order-zero bound.
The last assertion follows by passing to the limit on positive tests.  The
existence of a convergent subnet requires an appropriate locally convex
compactness theorem and is not asserted here without its hypotheses.
\end{proof}

\subsection{A compactness theorem that does hold}
\label{sec:ch14-compactness-positive}

The hypotheses alluded to at the end of the last proof can in fact be
supplied, provided one accepts nets rather than sequences.  We state the
result in the form in which it is used in
\Cref{chap:cld-bedford-taylor,part:bfj}.

\begin{theorem}
\label{thm:ch14-positive-compactness}
Let $X$ be locally holomorphically separated and paracompact, and let
$(T_i)_{i\in I}$ be a net of positive currents of bidimension $(p,p)$ on $X$.
Assume that for every relatively compact tropical chart $f:U\to\mathbf
G_m^r$, every compact $L\subset U$, and one nonnegative cutoff $\chi$ equal to
$1$ on $L$ and compactly supported in $U$, the masses
\[
 M_{U,L}:=\sup_i\bigl\langle T_i,\chi\,f_{\trop}^*\eta^{p}\bigr\rangle
\]
are finite, where $\eta=\dprimeddprime\|x\|^2$.  Then the net has a subnet
converging weakly to a positive current $T$ of bidimension $(p,p)$, and $T$
is of order $0$ relative to every chart datum.  If moreover every $T_i$ is
closed, then so is $T$.
\end{theorem}

\begin{proof}
\emph{Pointwise bounds.}  Fix $\alpha\in A_{X,c}^{p,p}$.  By definition of
$A_{X,c}^{p,p}$ there are finitely many relatively compact tropical charts
$f_j:U_j\to\mathbf G_m^{r_j}$ covering $\Supp\alpha$ and superforms
$\alpha_j$ with $\alpha|_{U_j}=f_{j,\trop}^*\alpha_j$.  Choose a subordinate
partition of unity $(\rho_j)$ with $\rho_j\ge0$ and, for each $j$, a compact
$L_j\supset\Supp\rho_j$ together with the cutoff $\chi_j$ of the hypothesis,
so that $\chi_j\ge\rho_j$.  The order-zero estimate proved inside
\Cref{prop:current-positivity-calculus} gives
\[
 \bigl|\langle T_i,\rho_j\alpha\rangle\bigr|
 \le c\,\bigl\langle T_i,\rho_jf_{j,\trop}^*\eta^p\bigr\rangle\,
   \|\alpha_j\|_{W_j,0},
\]
and $\langle T_i,\rho_jf_{j,\trop}^*\eta^p\rangle
\le\langle T_i,\chi_jf_{j,\trop}^*\eta^p\rangle\le M_{U_j,L_j}$, because
$(\chi_j-\rho_j)f_{j,\trop}^*\eta^p$ is a positive test form and $T_i$ is
positive.  Summing over $j$ gives a bound
$|\langle T_i,\alpha\rangle|\le C_\alpha$ independent of $i$.

\emph{Extraction.}  Consider the product space
$P=\prod_{\alpha}[-C_\alpha,C_\alpha]$, indexed by
$\alpha\in A_{X,c}^{p,p}$, with the product topology; it is compact by
Tychonoff's theorem.  The map $i\mapsto(\langle T_i,\alpha\rangle)_\alpha$
takes values in $P$, so the net has a convergent subnet
$(T_{i_\lambda})$; write $\langle T,\alpha\rangle$ for the limit of
$\langle T_{i_\lambda},\alpha\rangle$.  Convergence in $P$ is exactly
pointwise convergence on test forms, so $T_{i_\lambda}\to T$ in the weak
topology, provided $T$ is a current.

\emph{$T$ is a positive current.}  Linearity passes to the limit.  For a
positive $\alpha$ one has $\langle T_{i_\lambda},\alpha\rangle\ge0$, hence
$\langle T,\alpha\rangle\ge0$; and $\langle T,J\alpha\rangle=\langle
T,\alpha\rangle$ likewise.  Thus $T$ is a positive symmetric linear
functional, and \Cref{prop:ch14-automatic-continuity}\,(1) makes it a current
of order $0$ relative to every chart datum, with no further estimate.

\emph{Closedness.}  If $\dprime T_i=0$ for all $i$, then for a test form
$\beta$ of the appropriate bidegree
$\langle T,\dprime\beta\rangle=\lim\langle T_{i_\lambda},\dprime\beta\rangle
=\pm\lim\langle\dprime T_{i_\lambda},\beta\rangle=0$; the sign
\eqref{eq:ch14-dprime-sign} is the same for all $i$ because the bidegree is
fixed.  Hence $\dprime T=0$, and likewise for $\ddprime$.
\end{proof}

\begin{warning}{Subnets, not subsequences}
\Cref{thm:ch14-positive-compactness} produces a subnet.  A subsequence
cannot be expected: by \Cref{ex:ch14-star-unbounded-order} the space
$A_{X,c}^{p,p}$ is not normable, so the weak topology on bounded sets of
currents is not metrizable, and the usual sequential extraction is
unavailable.  This is one reason why the convention of
\Cref{chap:topological-tools}, that convergence statements are formulated for
nets unless a sequential reduction has been proved, is not a pedantry in this
part of the book.
\end{warning}

\subsection{What fails, relative to the complex case}
\label{sec:ch14-what-fails}

We now list precisely the parts of the classical package that are unavailable
here, with the example that obstructs each.  The list is short, and each item
has been proved above.

\paragraph{(F1) There is no global order, hence no normed duality.}
On a complex manifold, currents of order $\le m$ with support in a fixed
compact set form the dual of a Banach space, and Banach--Alaoglu applies.
Here, by \Cref{ex:ch14-star-unbounded-order}, a compact space---the Berkovich
projective line over a trivially valued field---carries a current of
bidimension $(0,0)$ whose order relative to a chart datum is unbounded as the
datum varies.  The space $A_{X,c}^{p,q}$ is an inductive limit over a
system of chart data admitting no cofinal sequence, and the resulting weak
topology on currents is not metrizable.  Consequently:
\begin{itemize}
\item ``bounded mass'' is not a meaningful hypothesis for a general current;
\item compactness statements must be net statements, as in
  \Cref{thm:ch14-positive-compactness};
\item there is no diagonal-extraction argument available.
\end{itemize}

\paragraph{(F2) Order zero does not mean measure coefficients.}
By \Cref{ex:ch14-star-not-measure} there is a current of bidimension $(0,0)$
on a compact Berkovich space which is of order $0$ relative to every chart
datum but is not induced by any signed Radon measure.  Over $\mathbf C$ this
cannot happen.  The failure is not exotic: it occurs on
$\mathbf P^{1,\An}_k$ for $k$ trivially valued.  Note that positivity
repairs it, by the last assertion of
\Cref{ex:ch14-forms-and-measures}; the phenomenon is confined to signed
currents.

\paragraph{(F3) There is no convolution, hence no canonical smoothing.}
On $\mathbf R^N$ every supercurrent can be regularized: for a nonnegative
$\rho\in C_c^\infty$ with $\int\rho=1$ and $\rho_\varepsilon(x)
=\varepsilon^{-N}\rho(x/\varepsilon)$, the convolution
$\rho_\varepsilon*T$ is the current of a \emph{smooth} form, commutes with
$\dprime$ and $\ddprime$, preserves each of the three positivity cones, and
converges weakly to $T$.  In particular every positive closed supercurrent on
$\mathbf R^N$ is a weak limit of smooth positive closed forms; this is
Lagerberg's regularization, recalled in \cite[\S1.6]{CLD25}.

Every step of that construction uses the group structure of $\mathbf R^N$
through three properties: translations act on the space, they act by
automorphisms of the sheaf of smooth forms, and the addition map is proper on
the relevant supports.  A Berkovich space has none of them.  There is no
translation action, the tropical charts have varying dimension, and---the
sharpest obstruction---the local tropical dimension is not locally constant,
so there is no ``direction'' in which to smear.  What replaces convolution is
chartwise smoothing followed by a comparison, and the comparison is exactly
what is missing: a chart is a map to $\mathbf R^r$ with no section.

Two positive substitutes are available and should not be confused with a
smoothing theorem.

\begin{proposition}
\label{prop:ch14-weak-density}
Let $X$ be boundaryless of pure dimension $n$, and assume the injectivity
statement of \Cref{rem:ch14-form-injectivity} on $X$, which holds for instance
whenever \cite[Corollary~7.3.9]{CLD25} applies.  Then the image of
$\omega\mapsto[\omega]$, from smooth $(p,q)$-forms to $D^{p,q}(X)$, is dense
in the weak topology.
\end{proposition}

\begin{proof}
By the Hahn--Banach theorem it suffices to show that a weakly continuous
linear functional on $D^{p,q}(X)$ vanishing on all $[\omega]$ vanishes
identically.  A weakly continuous functional on a space with the topology of
pointwise convergence on $A_{X,c}^{n-p,n-q}$ is evaluation against a test
form $\alpha$.  If $\int_X\omega\wedge\alpha=0$ for every smooth $\omega$,
then $[\alpha]=0$ as a current, whence $\alpha=0$ by the injectivity quoted
in \Cref{rem:ch14-form-injectivity}.
\end{proof}

\begin{warning}{Density is not regularization}
\Cref{prop:ch14-weak-density} says nothing about \emph{positive} or
\emph{closed} currents: the approximating smooth forms it produces are
neither.  A regularization theorem would have to produce, for a positive
closed $T$, smooth positive closed $\omega_\varepsilon$ with
$[\omega_\varepsilon]\to T$; and no such theorem is known in this
generality.  The distinction matters because every application in
\Cref{part:bfj,part:bj} that ``regularizes'' does so at the level of
\emph{potentials}, not of currents.
\end{warning}

Regularization of an arbitrary positive current by convolution is likewise
not intrinsic on a Berkovich space.  CLD define continuous psh functions by
a transfinite, sheafwise uniform-closure process.  Only the first-stage
``locally approachable'' class is, by definition, locally uniformly
approximable by smooth psh functions; stabilization at that stage is not
known.  Delta-forms and harmonic tropicalizations, treated in
\cref{chap:delta-forms},
offer different enlargements; they should not be described as a general
smoothing theorem for currents.

\paragraph{(F4) Removable sets are larger, not smaller.}
This is a failure only in the sense that intuition transported from
$\mathbf C$ misleads.  By \Cref{lem:ch14-extension-by-zero} and
\Cref{thm:ch14-support-dimension}, a closed set of tropical dimension
$<\max\{r,s\}$ is \emph{automatically} removable for currents of bidimension
$(r,s)$, with no growth or positivity hypothesis and with restriction an
isomorphism.  Over $\mathbf C$ the corresponding statement for a
$(1,1)$-current across a point on a surface requires positivity and the Skoda
or El Mir extension theory.  The reason for the discrepancy is that here the
\emph{test forms} already vanish on the small set, whereas over $\mathbf C$
they do not.  The two situations are therefore not comparable, and the
non-Archimedean statement must not be quoted as a strengthening of the
complex one.

\paragraph{(F5) Mass is not controlled by cohomology.}
Over a compact K\"ahler manifold the mass of a positive closed current is
determined by its cohomology class, so families in a fixed class are
automatically of bounded mass, and \Cref{thm:ch14-positive-compactness}
becomes unconditional.  Here no such mechanism exists in general: the
cohomology of the complex $(A_X^{\bullet,\bullet},\dprime,\ddprime)$ is not
known to be finite-dimensional in general, and the analogue of the Wirtinger
inequality would require a globally defined strongly positive
$(1,1)$-form, that is, a K\"ahler form, which a Berkovich space does not
carry canonically.  The mass hypothesis in
\Cref{thm:ch14-positive-compactness} therefore has to be verified by hand in
each application; in \Cref{part:bfj} it is supplied by an intersection
number on a model, and in \Cref{chap:cld-bedford-taylor} by the
Chern--Levine--Nirenberg inequalities.

\section{Dictionary, conventions, and unit tests}
\label{sec:ch14-dictionary}

\subsection{From this chapter to the rest of the book}
\label{sec:ch14-dict-forward}

\begin{center}
\small
\begin{tabular}{p{.30\textwidth}|p{.28\textwidth}|p{.32\textwidth}}
\textbf{Established here}&\textbf{Used as}&\textbf{Where}\\ \hline
test-form topology
  (\cref{def:ch14-chart-datum})
  &domain of weak convergence
  &\cref{prop:test-form-continuity-criterion},
   \cref{chap:cld-bedford-taylor}\\
relative order
  (\cref{def:ch14-relative-order})
  &substitute for distributional order
  &\cref{sec:ch14-what-fails}\\
sheaf property
  (\cref{prop:ch14-currents-sheaf})
  &locality of curvature
  &\cref{chap:cld-poincare-lelong}\\
extension by zero
  (\cref{lem:ch14-extension-by-zero})
  &removability across thin sets
  &\cref{chap:cld-bedford-taylor},
   \cref{chap:cld-metrized-bundles}\\
sign of $\dprime$ on currents
  (\cref{lem:ch14-sign-forced})
  &Poincar\'e--Lelong normalization
  &\eqref{eq:appJ-cld-pl},
   \cref{chap:cld-poincare-lelong}\\
boundary current
  (\cref{prop:ch14-boundary-current})
  &Stokes with boundary
  &\cref{chap:cld-forms},
   \cref{sec:appJ-currents}\\
positivity hierarchy for currents
  (\cref{prop:ch14-form-current-cones})
  &positivity of curvature currents
  &\cref{chap:cld-metrized-bundles}\\
automatic continuity
  (\cref{prop:ch14-automatic-continuity})
  &construction of $\delta_Y$, $\delta_\Sigma$, $u_\omega$
  &\cref{sec:ch14-cycles},
   \cref{sec:ch14-fibre-integration}\\
proper pushforward
  (\cref{prop:current-pushforward-formal})
  &projection formula, degrees
  &\cref{cor:ch11-projection-formula},
   \cref{chap:cld-metrized-bundles}\\
cycle current
  (\cref{thm:current-fundamental-cycle})
  &right-hand side of Poincar\'e--Lelong
  &\cref{chap:cld-poincare-lelong}\\
fibre integration
  (\cref{thm:cld-fibre-integration})
  &pushforward of Monge--Amp\`ere measures
  &\cref{thm:cld-fibre-projection},
   \cref{chap:cld-metrized-bundles}\\
support theorem
  (\cref{thm:ch14-cld-support-corrected-packet})
  &polar negligibility
  &\cref{chap:cld-bedford-taylor}\\
weak compactness
  (\cref{thm:ch14-positive-compactness})
  &limits of Monge--Amp\`ere measures
  &\cref{part:bfj},
   \cref{chap:cld-bedford-taylor}\\
weakly null directions
  (\cref{ex:ch14-weakly-null})
  &hypothesis in convex products
  &\cref{thm:ch14-convex-bt-imported},
   \cref{sec:appD-supercurrents}\\
strongly homogeneous currents
  (\cref{prop:ch14-strongly-homogeneous})
  &mixed volumes, Alexandrov--Fenchel
  &\cref{thm:appD-support-current-mixed-volume},
   \cref{thm:appD-af}
\end{tabular}
\end{center}

\subsection{The convention dictionary}
\label{sec:ch14-convention-dictionary}

Four bodies of literature meet in this chapter---the Lagerberg--CLD calculus,
Berndtsson's 2019 supercurrent paper, his 2026 convex paper, and classical
complex pluripotential theory---and they do not agree on signs, on the
meaning of the word ``positive'', or on where a factorial goes.  The table
below is the promised reconciliation.  Every row is a definition, not a
theorem; the diagnostic column is the computation that detects a mismatch.
The binding normalizations of this book are those of
\Cref{sec:appJ-cld,sec:appJ-currents}, and nothing in this table overrides
them.

\begin{center}
\small
\begin{tabular}{@{}p{.20\textwidth}p{.25\textwidth}p{.25\textwidth}p{.24\textwidth}@{}}
\toprule
Datum & This book (Lagerberg--CLD) & Elsewhere & Diagnostic \\
\midrule
Bigrading
 & $\dprime x_i$, $\ddprime x_i$; bidegree $(p,q)$
 & $\rd x_i$, $\rd\xi_i$ in \cite{Ber19Super,Ber26Convex}
 & identify $\rd\xi_i\leftrightarrow\ddprime x_i$ \\
Involution
 & $J\dprime x_i=\ddprime x_i$, $J\ddprime x_i=\dprime x_i$
 & $J\rd x_i=-\rd\xi_i$, $J\rd\xi_i=\rd x_i$ in
   \cite[\S2]{Ber26Convex}; opposite sign in \cite{Ber19Super}
 & test $J^2$: ours is an involution \\
Second-order operator
 & $\dprimeddprime u=\sum u_{ij}\dprime x_i\wedge\ddprime x_j$
 & $\rd\rd^{\#}u$ with $\rd^{\#}=J^{-1}\rd J$
 & convex $u$ must give a positive form,
   \eqref{eq:appJ-real-hessian} \\
Top orientation
 & $\tau_n=\bigwedge_i(\dprime x_i\wedge\ddprime x_i)$, interlaced
 & grouped $\dprime x_{[n]}\wedge\ddprime x_{[n]}$; Berezin sign
 & the constant is $c_n$, \eqref{eq:appJ-cp} \\
``Positive''
 & three cones $\mathrm{SP}\subset\mathrm P\subset\mathrm{WP}$,
   \cref{def:ch09-cones}
 & complex usage: ``positive'' often means our weakly positive
 & only bidegrees $2\le p\le n-2$ distinguish them \\
``Strong''
 & \emph{strongly positive} $=$ in $\mathrm{SP}$
 & \emph{strong} in \cite{Ber26Convex} $=$ in the linear span of
   $\mathrm{SP}$
 & a weakly null form is neither \\
Positivity for currents
 & tested against the \emph{opposite} cone,
   \cref{prop:current-positivity-calculus}
 & sometimes tested against the same cone
 & \cref{prop:ch14-form-current-cones}: $[\omega]$ and $\omega$ agree \\
Current differential
 & $\dprime=(-1)^{p+q+1}b'$ on $D^{p,q}$
 & naked transpose $b'$
 & $\dprime[\omega]=[\dprime\omega]$,
   \cref{lem:ch14-sign-forced} \\
Boundary
 & $\dprime[\omega]=[\dprime\omega]-[\omega]^\partial$
 & no boundary term when $\partial X=\emptyset$
 & \cref{prop:ch14-boundary-current} \\
Bidegree of a cycle
 & $\delta_Z$ of codimension-$p$ cycle has bidegree $(p,p)$
 & bidimension $(n-p,n-p)$
 & \cref{sec:appJ-currents} \\
Monge--Amp\`ere factor
 & $\tfrac1{n!}(\dprimeddprime u)^n=\det(D^2u)\,\rd x$
 & $(\rd\rd^\#u)^n$ without $1/n!$
 & \cref{thm:appD-alexandrov-measure},
   \cref{prop:ch09-ma-coordinates} \\
Integration current of $M$
 & calibrated: lattice measure, $\delta_M$
 & metric: $[M]_{\mathrm s}$ of \cite{Ber19Super}
 & ratio $\operatorname{covol}_g(N_M)$,
   \cref{prop:ch14-berndtsson-vs-tropical} \\
Mean curvature
 & $\vec H=\sum_i\mathrm{II}(e_i,e_i)$, unnormalized
 & sometimes divided by $m$
 & sphere of radius $R$: $|\vec H|=(N-1)/R$ \\
Complex comparison
 & $\dprimeddprime\log|f|=\delta_{\operatorname{div}f}$
 & $\ddc\log|f|^2=[\operatorname{div}f]$
 & factor $2$ in the weight,
   \cref{sec:appJ-complex-ddc} \\
\bottomrule
\end{tabular}
\end{center}

Three rows deserve a sentence of commentary, because they are the ones that
have caused published confusion.

First, the word \emph{strong}.  In \cite{Ber26Convex} a form is strong when
it lies in the linear span of the strongly positive cone.  By
\Cref{cor:ch14-sp-not-spanning} that span is proper in bidegree $(2,2)$ on
$\mathbf R^4$, so ``strong'' is a genuine restriction and not a synonym for
``strongly positive'': a strongly positive form is strong, but a strong form
need not be positive in any sense.  The hypothesis in
\Cref{thm:ch14-convex-bt-imported} is the linear-algebraic one.

Second, the direction of the positivity hierarchy.  For forms the smallest
cone is $\mathrm{SP}$; for currents the smallest class is also called
strongly positive, but it is defined by testing against the \emph{largest}
cone of forms.  \Cref{thm:ch09-positivity-duality} is what makes the two
usages compatible, and \Cref{prop:ch14-form-current-cones} is the statement
that the compatibility is exact on currents defined by forms.

Third, the sign of $J$.  The two Berndtsson papers differ from each other,
and both differ from the Lagerberg--CLD convention used here.  Since
$\dprimeddprime$ is defined through $J$, no identity involving $J$ may be
transported without recomputation; identities involving only
$\dprimeddprime$ on \emph{functions} may, because all three conventions give
the Hessian in the ordered basis.  This is the content of
\Cref{rem:ch14-convention-firewall}.

\subsection{Unit tests}
\label{sec:ch14-unit-tests}

The chapter contains the following complete checks.  Together they exercise
every convention in the preceding table at least once.

\begin{center}
\small
\begin{tabular}{p{.34\textwidth}|p{.28\textwidth}|p{.30\textwidth}}
\textbf{Test}&\textbf{Detects}&\textbf{Where}\\ \hline
$\dprimeddprime\max\{0,t\}=\delta_0$
  &sign of the current differential
  &\cref{ex:ch14-convex-corner-restated}\\
model function on a star
  &vertex atom is minus the outgoing flux
  &\cref{ex:ch14-line-model-function}\\
$k$-th derivative on the $k$-th branch
  &no global order
  &\cref{ex:ch14-star-unbounded-order}\\
nonsummable $(c_b)$ on a star
  &order zero is not measure
  &\cref{ex:ch14-star-not-measure}\\
$\omega=-(\dprime x_{12}+\dprime x_{34})
   \wedge(\ddprime x_{12}+\ddprime x_{34})$
  &positive $\neq$ strongly positive
  &\cref{ex:ch14-positive-not-strong-current}\\
Pl\"ucker form $\widehat P$
  &existence of weakly null directions
  &\cref{ex:ch14-weakly-null}\\
$T$ and $-T$ positive
  &saliency of the positive cone only
  &\cref{prop:ch14-positive-and-negative}\\
$u_\varepsilon(x)=\chi(x/\varepsilon)$ at $0$
  &nonexistence of pullback
  &\cref{prop:ch14-no-pullback}\\
Tate curve skeleton
  &mass equals $v_K(q)$
  &\cref{ex:ch14-skeleton-current}\\
projection $\mathbf G_m^2\to\mathbf G_m$
  &unimodular case, no index
  &\cref{ex:ch14-fibre-projection-numbers}\\
doubling $S=T_1^2$
  &calibrated fibre degree $2$; total mass $4$
  &\cref{ex:ch14-fibre-degree-numbers}\\
$M=\mathbf R(1,1)\subset\mathbf R^2$
  &metric versus lattice normalization
  &\cref{ex:ch14-diagonal-line}\\
$g\mapsto\lambda^2g$
  &$[M]_{\mathrm s}\mapsto\lambda^{-m}[M]_{\mathrm s}$
  &\cref{prop:ch14-berndtsson-metric-dependence}\\
$h$ homogeneous of degree $1$
  &$\iota'_E\dprimeddprime h=0$
  &\cref{prop:ch14-strongly-homogeneous}
\end{tabular}
\end{center}

Five points are worth memorizing on leaving the chapter.
\begin{enumerate}
\item A current has an order relative to a chart datum, never in general a
  global one; every compactness statement must be read accordingly.
\item Duality reverses the two extreme positivity cones for the \emph{test}
  forms, but not for the objects: a current defined by a positive form is
  positive.
\item The only sign created by the current formalism is
  $(-1)^{p+q+1}$, and it is forced by
  $\dprime[\omega]=[\dprime\omega]$, not chosen.
\item Positivity buys continuity and order zero; it does not buy
  regularization.
\item ``Integration current of $M$'' means two different things, differing by
  the covolume of the tangent lattice, and only one of them is metric-free.
\end{enumerate}

\section*{Source notes}

The current topology, sheaf property, positivity, and differential calculus
are reconstructed from \cite[Chapter~7]{CLD25}; the analytic cycle and fibre
integration results are from \cite[Chapter~16]{CLD25}.  Our insistence on
proper support, dense flat locus, and the nontrivially valued hypothesis in
the arbitrary-current clause of the support theorem follows the exact
statements of that source.

\paragraph{Provenance of the individual statements.}
The definition of a current on a tropical space, the description of
continuity relative to a covering chart datum, and the observation that the
resulting orders need not be bounded are
\cite[\S7.1]{CLD25}; our
\Cref{ex:ch14-star-unbounded-order} makes that observation explicit on the
star of \cite[\S6.7.3]{CLD25}, which is the Berkovich projective line over a
trivially valued algebraically closed field, and
\Cref{ex:ch14-star-not-measure} is the example of
\cite[7.2.14]{CLD25} with its summability criterion proved in both
directions.  The sheaf property and fineness are
\cite[Proposition~7.1.7 and Corollary~7.1.8]{CLD25}; the extension of a
current to forms whose support meets $\Supp T$ compactly is
\cite[Lemma~7.1.10]{CLD25}; extension by zero across a set of small tropical
dimension is \cite[Lemma~7.1.13]{CLD25}.  The three positivity cones for
currents, their sheaf property, and the wedge rules are
\cite[\S7.2]{CLD25}, resting on the superform statements of
\cite[\S1.2]{CLD25} which we proved in
\Cref{chap:superforms}.  The identification of positive $(0,0)$-currents with
positive Radon measures is \cite[Proposition~7.2.8]{CLD25}, and the automatic
continuity of positive functionals is \cite[Proposition~7.2.9]{CLD25}; on the
latter see \Cref{rem:ch14-third-case}, where we explain why we prove and use
only two of the three cases.  Integration currents and their positivity are
\cite[\S7.3]{CLD25}; the differential calculus, including the boundary term
of \Cref{prop:ch14-boundary-current} and the sign
$(-1)^{p+q+1}$, is \cite[\S7.7]{CLD25}.  The cycle statement
\Cref{thm:current-fundamental-cycle} is
\cite[Proposition~16.6.3]{CLD25}, fibre integration is
\cite[Theorem~16.5.4 and Corollary~16.5.5]{CLD25}, the projection formula for
Bedford--Taylor products is \cite[Theorem~16.5.7]{CLD25}.  The
arbitrary-current support assertion is \cite[Theorem~16.8.5]{CLD25}.  The
trivially valued Bedford--Taylor assertion uses
\cite[Proposition~16.8.2 and Proposition~8.7.5]{CLD25} and
\cite[Lemma~16.8.7 and the proof of Corollary~16.8.6]{CLD25}, with the
exponent correction explained in
\cref{rem:ch14-cld-support-corrected-index}.

\paragraph{The Euclidean support input.}
The Euclidean vanishing statement used inside
\cite[Theorem~16.8.5]{CLD25}, and recalled in
\Cref{rem:ch14-support-route}, is
\cite[Corollary~4.1]{Lag12Super}; it is Corollary~3.5 in
arXiv:1008.2856v1.  That paper is also the origin of the
superform calculus of \Cref{chap:superforms} and of the convolution
regularization discussed in \Cref{sec:ch14-what-fails}; the reader should
consult \cite[\S\S1.2 and 1.6]{CLD25} for the statements in the form used
above.

\paragraph{The two Berndtsson constructions.}
The metric supercurrent of a real submanifold, its positivity, the
identification of its traced differential with the first variation of area,
and the mean-curvature flow variation formula are \cite[\S\S4--5]{Ber19Super};
we use the corrected 2019 version, in which the argument of Section~8 of the
first arXiv version was replaced while the main results were unchanged.
\Cref{lem:ch14-berndtsson-pairing},
\Cref{prop:ch14-berndtsson-metric-dependence},
\Cref{prop:ch14-berndtsson-vs-tropical}, and
\Cref{ex:ch14-diagonal-line} are not in that paper; they are the
translation into the calibrated language of this book, and the covolume
factor they isolate appears to be the cleanest way to state that
$[M]_{\mathrm s}$ and $\delta_M$ are different objects.
\Cref{prop:ch14-first-variation} is the classical first-variation formula and
is included so that the imported statement has a visible content.  The
convex theory---weakly null directions, the notion of a strong form, the
Bedford--Taylor product for convex potentials, strongly homogeneous
supercurrents, and the Lorentzian form behind Alexandrov--Fenchel---is
\cite[\S\S3--6 and 10]{Ber26Convex}.  The mixed-volume identity is proved in
\Cref{thm:appD-support-current-mixed-volume}, whose only analytic input is
\Cref{thm:ch14-convex-bt-imported}; the convex-analytic background is
\cite{Roc70} and \cite{Sch14}.

\paragraph{Departures from the sources.}
Three points of presentation differ deliberately from
\cite{CLD25}.  First, we prove automatic continuity only for functionals
nonnegative on positive, or on weakly positive, forms, because
\Cref{cor:ch14-sp-not-spanning} shows that the domination argument does not
reach the third case; this costs nothing, since every current constructed in
this book by positivity is of the first two kinds.  Second, we make the
failure of a global order into a numbered example rather than a remark,
because it is the single fact from which the whole of
\Cref{sec:ch14-what-fails} follows.  Third, we state and prove the weak
compactness theorem \Cref{thm:ch14-positive-compactness} in net form, using
Tychonoff, rather than leaving the hypotheses unstated; the classical
complex counterpart is in \cite[Chapter~III]{Dem12}, and the comparison
there is the point of \Cref{sec:ch14-what-fails}.
